% Options for packages loaded elsewhere
\PassOptionsToPackage{unicode}{hyperref}
\PassOptionsToPackage{hyphens}{url}
\documentclass[
]{article}
\usepackage{xcolor}
\usepackage[margin=1in]{geometry}
\usepackage{amsmath,amssymb}
\setcounter{secnumdepth}{-\maxdimen} % remove section numbering
\usepackage{iftex}
\ifPDFTeX
  \usepackage[T1]{fontenc}
  \usepackage[utf8]{inputenc}
  \usepackage{textcomp} % provide euro and other symbols
\else % if luatex or xetex
  \usepackage{unicode-math} % this also loads fontspec
  \defaultfontfeatures{Scale=MatchLowercase}
  \defaultfontfeatures[\rmfamily]{Ligatures=TeX,Scale=1}
\fi
\usepackage{lmodern}
\ifPDFTeX\else
  % xetex/luatex font selection
  \setmainfont[]{Cambria}
  \setmonofont[]{Consolas}
\fi
% Use upquote if available, for straight quotes in verbatim environments
\IfFileExists{upquote.sty}{\usepackage{upquote}}{}
\IfFileExists{microtype.sty}{% use microtype if available
  \usepackage[]{microtype}
  \UseMicrotypeSet[protrusion]{basicmath} % disable protrusion for tt fonts
}{}
\makeatletter
\@ifundefined{KOMAClassName}{% if non-KOMA class
  \IfFileExists{parskip.sty}{%
    \usepackage{parskip}
  }{% else
    \setlength{\parindent}{0pt}
    \setlength{\parskip}{6pt plus 2pt minus 1pt}}
}{% if KOMA class
  \KOMAoptions{parskip=half}}
\makeatother
\usepackage{longtable,booktabs,array}
\newcounter{none} % for unnumbered tables
\usepackage{calc} % for calculating minipage widths
% Correct order of tables after \paragraph or \subparagraph
\usepackage{etoolbox}
\makeatletter
\patchcmd\longtable{\par}{\if@noskipsec\mbox{}\fi\par}{}{}
\makeatother
% Allow footnotes in longtable head/foot
\IfFileExists{footnotehyper.sty}{\usepackage{footnotehyper}}{\usepackage{footnote}}
\makesavenoteenv{longtable}
\usepackage{graphicx}
\makeatletter
\newsavebox\pandoc@box
\newcommand*\pandocbounded[1]{% scales image to fit in text height/width
  \sbox\pandoc@box{#1}%
  \Gscale@div\@tempa{\textheight}{\dimexpr\ht\pandoc@box+\dp\pandoc@box\relax}%
  \Gscale@div\@tempb{\linewidth}{\wd\pandoc@box}%
  \ifdim\@tempb\p@<\@tempa\p@\let\@tempa\@tempb\fi% select the smaller of both
  \ifdim\@tempa\p@<\p@\scalebox{\@tempa}{\usebox\pandoc@box}%
  \else\usebox{\pandoc@box}%
  \fi%
}
% Set default figure placement to htbp
\def\fps@figure{htbp}
\makeatother
\setlength{\emergencystretch}{3em} % prevent overfull lines
\providecommand{\tightlist}{%
  \setlength{\itemsep}{0pt}\setlength{\parskip}{0pt}}
\AtBeginDocument{%
  \author{%
    Mikhail L. Arbuzov\\
    Independent Researcher\\
    \texttt{mike.arbuzov54@gmail.com}
    \and
    Sisong Bei\\
    Independent Researcher\\
    \texttt{qurining@gmail.com}
    \and
    Ziwei Dong\\
    Independent Researcher\\
    \texttt{ziwei.dong@alumni.emory.edu}
    \and
    Dmitri Kalaev\\
    Independent Researcher\\
    \texttt{kalaevdr@gmail.com}
    \and
    Alexey Shvets\\
    Palo Alto Networks\\
    \texttt{ashvets@paloaltonetworks.com}
  }%
}
\usepackage{bookmark}
\IfFileExists{xurl.sty}{\usepackage{xurl}}{} % add URL line breaks if available
\urlstyle{same}
\hypersetup{
  pdftitle={The Architecture of Errors: Logarithmic Mode Discovery and Polylogarithmic Intervention Budgets for Long-Context LLM Reliability},
  hidelinks,
  pdfcreator={LaTeX via pandoc}}

\title{The Architecture of Errors: Logarithmic Mode Discovery and
Polylogarithmic Intervention Budgets for Long-Context LLM Reliability}
\author{}
\date{}

\begin{document}
\maketitle

\textbf{Abstract}

Universal LLM reliability is not a finite-library problem: across all
possible tasks, tools, schemas, knowledge sources, and evaluator
expectations, new intervention-distinguishable failure modes can appear
without bound, so no finite intervention dictionary can guarantee
bounded residual error for every intervention-distinguishable mode of an
unbounded domain. But deployed systems do not operate over the whole
universe; they operate inside operationally bounded patches with
recurring tasks, schemas, tools, and evaluator expectations. Within such
patches, empirical evidence suggests failures are sparse, repetitive,
and concentrated in a small recurring catalogue, so reliability becomes
a local catalogue-discovery and intervention-coverage problem rather
than an exponential token-length problem. We formalize this transition
with a worst-case-mode-wise negative result showing that no finite
intervention dictionary covers every intervention-distinguishable
failure mode of an unbounded domain, an inverse-discovery corollary
showing that the logarithmic upper envelope cannot deliver linearly more
distinct tail modes without exponentially more observed hard-failure
events, and a positive patch-local result showing that under log
active-mode exposure and head-heavy coverage, a sufficient
per-hard-decision intervention budget grows polylogarithmically in
sequence length and becomes domain-constant once the patch catalogue
saturates. The framework relocates rather than dissolves long-context
difficulty: where the number of hard decisions itself grows with task
length, reliability remains hard; the contribution is to identify the
on-axis intervention rather than to make those regimes easy.

\begin{center}\rule{0.5\linewidth}{0.5pt}\end{center}

\subsection{1. Introduction}\label{introduction}

The standard worry about long-context autoregressive generation is
exponential. If every token has independent error probability \(e\),
then after \(n\) tokens the chance of a fully correct output is
\((1-e)^n\), which collapses to zero for any nontrivial \(e\) and large
enough \(n\) (LeCun, 2023; Dziri et al., 2023). In our previous paper
(Arbuzov et al., 2025) we argued that this worry is misplaced because
errors are not distributed uniformly across tokens. Empirical work (Fang
et al., 2025) shows that only 5--10\% of tokens are ``key'' ---
genuinely dependent on long-range context --- while the remaining
majority become nearly deterministic once enough context accumulates.
Replacing \((1-e)^n\) with the two-rate model

\[
P(\text{correct}) = (1-e_{\text{key}})^{k} \cdot (1-e_{\text{non}})^{n-k}, \qquad k \ll n,
\]

with \(k\) scaling sublinearly in \(n\), recovers the observed
long-context coherence of modern LLMs and converts the reliability
question from ``how does \(n\) grow?'' to ``how does \(k\) grow?''.

\subsubsection{1.1 From sparse tokens to recurring
patterns}\label{from-sparse-tokens-to-recurring-patterns}

This paper takes the next step. Sparsity tells us \emph{where} errors
live. The natural follow-up question is \emph{what} they are. Recent
failure-mode atlases supply a striking empirical answer: errors are not
only sparse, they are also \emph{repetitive}. The ErrorAtlas project
(Ashury-Tahan et al., 2026), covering 83 models across 35 datasets,
organises observed failures into 17 named categories sorted by
prevalence, with a long-tailed distribution heavily concentrated in the
head. In code, two error types (\texttt{AssertionError} and
\texttt{NameError}) cover 86.35\% of failures on HumanEval across 14
LLMs, a result replicated across 23 models on HumanEval Pro and MBPP Pro
(Wen et al., 2024). In math, MWPES-300K (Sun et al., 2025) categorises
304,865 errors from 15 LLMs across four math word-problem datasets and
reports that dataset characteristics shape error patterns
systematically, with a small number of dominant categories per model.
The same picture appears in multi-hop QA (Zhang et al., 2026), agentic
tool use (Cemri et al., 2025), and retrieval-augmented generation (Wood
\& Forbes, 2024).

These observations suggest a third layer in the reliability
architecture, complementing key-token sparsity (Layer 1) and the
within-key stratified-manifold structure (Layer 2 in Arbuzov et al.,
2025). Within the small set of key tokens, only a fraction \(\beta\)
produce \emph{hard} failures, and those hard failures cluster into a
finite catalogue of recurring modes whose size grows much more slowly
than the number of observed events.

\subsubsection{1.2 The intervention
budget}\label{the-intervention-budget}

If failure modes form a small, slowly-growing catalogue, then targeted
interventions become a viable engineering strategy. Each recurring mode
admits a specific countermeasure: arithmetic errors fall to a Python
interpreter (Gao et al., 2022; Chen et al., 2022; Gou et al., 2023);
format violations to constrained decoding (Suresh et al., 2025; Wang D.
Y.-B. et al., 2025); hallucinations to retrieval augmentation, with one
paper reporting 100\% elimination on RAGTruth (Wood \& Forbes, 2024;
95\% CI 91--100\%, benchmark-conditional --- not a by-construction
guarantee for RAG in general); over-refusals to preference optimisation
(Karaman et al., 2024); tool-call errors to structured uncertainty (Suri
et al., 2025). We collect twelve such categories with quantitative ``one
cluster, one intervention'' results in Section 4.2.

The natural question becomes: how many interventions are enough? If new
failure modes follow canonical Heaps-style growth (Manning et al.,
2008), the optimistic doubly-logarithmic rate fails, but the
intervention budget can still remain polylogarithmic when the number of
hard decisions grows sublinearly. The practical question is therefore
the exponent and constants, not whether the program is hopeless. The
data also suggest the more optimistic case: error-taxonomy work plateaus
quickly, with most domains stabilising at 10--50 recurring modes even at
large corpus sizes. We treat this as an \emph{empirical postulate of
logarithmic mode discovery} (Section 3.2) rather than as a derived
consequence --- Zipfian rank-frequency does \emph{not} imply logarithmic
cluster growth in the strict mathematical sense; Heaps' law is the
correct type--token consequence of Zipf-distributed events, and it is
power-law, not logarithmic. Under our postulate, a sufficient
per-hard-decision intervention budget scales polylogarithmically in
sequence length under log active-mode exposure (Proposition 2, Section
3.4).

\subsubsection{1.3 Contributions}\label{contributions}

The paper's central contribution is a shift in the reliability object:
universal LLM reliability is not a finite-library problem, but
patch-local reliability can be treated as catalogue discovery and
intervention coverage. We formalise this transition with two
propositions and one corollary. Proposition 1 gives the
worst-case-mode-wise negative result: if an unbounded domain keeps
producing intervention-distinguishable failures, no finite intervention
dictionary can guarantee bounded residual error for every
intervention-distinguishable mode of the domain. Corollary 1 gives the
inverse economic implication of the logarithmic upper envelope: the
envelope cannot deliver linearly more distinct tail modes without
exponentially more observed hard-failure events, so open-domain tail
discovery has rapidly diminishing returns. Proposition 2 gives the
corresponding sufficient, model-implied positive result inside a fixed
deployment patch: under log active-mode exposure and head-heavy
coverage, a sufficient per-hard-decision intervention budget satisfies
\(m \geq \lceil |C_{\text{eff}}|^{1 - \varepsilon / e_{\text{hard}}} \rceil\),
with a doubly-logarithmic sequence-length rate as the optimistic pre-cap
special case and a domain-constant budget once the patch catalogue
saturates. The sequence-level analogue is strictly tighter and
approaches full-catalogue coverage as \(k\) grows; we surface that
asymmetry rather than bury it.

Around this transition the paper does four supporting things. A
three-layer framework --- sparsity (\(\alpha\)), hard-token
stratification (\(\beta\)), and patch-local mode catalogue (\(|C_D|\))
--- refines our previous two-rate model and separates \emph{where}
errors occur, \emph{what} recurring forms they take, and \emph{which}
capability interventions address them. Logarithmic mode discovery is
stated as an empirical postulate, not a theorem, calibrated against
ErrorAtlas, HumanEval, and MWPES with \(\sigma \in [0.87, 1.85]\) across
anchors and \(\sigma \approx 1.85\) carried as a conservative planning
value. Section 4 synthesises evidence for failure clustering,
cluster-selective interventions, and sublinear length scaling, drawing
on approximately sixty prior published results, including the six-axis
capability-elimination harvest of 28 quantitatively-anchored citations
stratified into Patterns A/B/C (by-construction, strong empirical with
class-shift, moderate with residual shift). And Section 6 re-audits the
most-cited steep-decay counter-evidence (Dziri et al., 2023; Kuratov et
al., 2024; Kwa et al., 2025; Wan et al., 2025), showing that each decays
primarily over task-structure variables --- compositional graph size,
fact count, log-time horizon, capacity threshold, or evidence scope ---
rather than raw token length.

\begin{center}\rule{0.5\linewidth}{0.5pt}\end{center}

\subsection{2. Related Work}\label{related-work}

\subsubsection{2.1 Failure-mode
taxonomies}\label{failure-mode-taxonomies}

The systematic cataloguing of LLM failures has reached maturity across
several domains in the last two years. ErrorAtlas (Ashury-Tahan et al.,
2026) provides the largest cross-domain taxonomy to date: 83 models
evaluated on 35 datasets yield approximately 7,000 sampled failures
organised into 17 named categories sorted by prevalence (the source's
Table 1), with a long-tailed, head-concentrated distribution. MWPES-300K
(Sun et al., 2025) takes the math domain deeper, with 304,865
categorised errors from 15 LLMs across four math word-problem datasets
and a strict per-model decomposition that shows the same top-4
categories dominating across models --- a striking dataset-invariance
result that we use heavily in Section 4.1. For code, Wen et al.~(2024)
categorise 12,837 errors from 14 LLMs on HumanEval, then replicate the
dominant-category hierarchy on 23 models across HumanEval Pro and MBPP
Pro. The RFMDataset (Guo et al., 2025) explicitly tests cross-model
invariance and reports ``strikingly similar failure mode distributions''
across 10 advanced reasoning models (three open-source and seven
proprietary).

Beyond general-purpose evaluations, domain-specific taxonomies exist for
multi-hop QA (the Weakest Link Law authors (2026) show that multi-hop
performance collapses to the least-visible-evidence bucket on MuSiQue
and NeoQA), for multi-agent systems (Cemri et al., 2025), and for
tool-augmented agents (Yao et al., 2024). The convergent finding across
all these is the same: at any given corpus size, the number of distinct
named modes is small (typically 8--20), and the cumulative coverage of
the top modes is high. This is the empirical phenomenon our theoretical
framework aims to model.

\subsubsection{2.2 Targeted interventions}\label{targeted-interventions}

Each named failure mode in these taxonomies has, in some cases
independently, become the subject of focused intervention research.
Arithmetic is the canonical example. PAL (Gao et al., 2022) delegates
execution to a Python interpreter and lifts GSM-Hard accuracy from
20.1\% (chain-of-thought) to 61.5\% --- a 41.4-point gain whose
residuals are problem-comprehension errors rather than arithmetic ones.
PoT (Chen et al., 2022) and ToRA (Gou et al., 2023) extend the same
delegation to broader math and financial reasoning. Format violations
admit something stronger: constrained decoding zeroes out invalid-token
probability at every decoding step, which is mathematically rather than
statistically a fix --- DINGO (Suresh et al., 2025) reports up to +68
percentage points on combined symbolic-math and JSON-generation
benchmarks, and SLOT (Wang D. Y.-B. et al., 2025) achieves 99.5\% schema
accuracy while preserving 94\% content similarity. Code-logic errors
yield to execution feedback (Reflexion: 80\% → 91\% HumanEval pass@1;
AgentCoder pushing the same benchmark to 96.3\%). Step-level reasoning
errors yield to process reward models: Math-Shepherd (Wang P. et al.,
2023) raises Mistral-7B MATH accuracy from 28.6\% to 33.0\% under
step-by-step PPO, then to 43.5\% when the same reward model is used for
verification reranking on top; Let's Verify Step by Step (Lightman et
al., 2023) reaches 78\% on the MATH test set under process supervision.
Retrieval-augmented generation is the most studied of all, with Acurai
(Wood \& Forbes, 2024) reporting 100\% hallucination elimination on
RAGTruth --- a benchmark-conditional empirical observation, not a
property of RAG in general --- and corroborating clinical and legal
results elsewhere. Two more clusters round out the picture. POROver
(Karaman et al., 2024) lifts the not-overrefused rate from 57.6\% to
82.1\% while preserving 97.9\% safety; SAGE-Agent (Suri et al., 2025)
raises When2Call accuracy from 36.5\% to 65.2\%.

Two recurring observations from this literature inform our theoretical
framework. First, each intervention is \emph{cluster-specific}: the
residual errors after intervention belong to a different named cluster,
not to the targeted one. Reflexion residuals are
specification-misinterpretation errors, not execution errors;
constrained-decoding residuals are semantic errors within valid
structure, not format violations. This is the empirical signature of the
cluster-orthogonality property that anchors Section 4.2 --- though, more
precisely, the empirical evidence supports \emph{selectivity}
(post-intervention residuals shift outside the targeted class) rather
than literal orthogonality (perfect independence of failure classes);
the framework only requires the former. Second, \emph{additivity is
approximate but not perfect}. The Six Sigma Agent framework reports
compound reliability gains of \(14{,}700\times\) when interventions
target orthogonal failure modes (Patel et al., 2026), but Le (2026)
shows that schema-level and prompt-level instructions interact
non-additively when they share a prompt channel. We handle this caveat
directly in Section 5.

\subsubsection{2.3 Length-decay benchmarks and their
interpretation}\label{length-decay-benchmarks-and-their-interpretation}

A parallel literature measures the \emph{shape} of the reliability decay
curve as context length grows. Two camps exist. The first reports decay
that is much milder than exponential. Loong (Wang M. et al., 2024) finds
GPT-4o on Chain-of-Reasoning tasks declines from 81.6\% (10--50K tokens)
to 32.9\% (\textgreater200K) --- a relative drop of about 60\% across
two doublings, consistent with log-linear decay over \(\log L\) across
the measured bins. We report the decay shape qualitatively rather than
committing to fitted coefficients; the bin-representative convention
materially shifts the slope, while the qualitative claim is unaffected.
GSM-\(\infty\) (Zhou et al., 2025) fits a sigmoid (saturating) decay
with \(R^2 > 0.98\) and shows that exponentially more inference compute
yields only linear gains in AUC --- a direct signature of clustered, not
independent, errors. RULER (Hsieh et al., 2024) finds threshold-like
behaviour where simple retrieval holds to 32K while complex reasoning
collapses earlier, inconsistent with smooth \((1-\varepsilon)^n\).
Anchor-based LLMs (Pang et al., 2024) achieve 99\% reduction in cached
K/V state with only ``minor accuracy compromise,'' which implies that
the effective number of distinct attention-key vectors the model needs
to retain is far smaller than \(n\) --- an architectural correlate of
the sparse-decision claim, though K/V compression is not literal token
removal.

The second camp reports ``rapid decay'' findings that are sometimes
cited as evidence for exponential compounding. Faith and Fate (Dziri et
al., 2023) shows multi-digit multiplication accuracy collapsing from
59\% at 3-digit to 4\% at 4-digit (zero-shot GPT-4) and argues
theoretically that ``the probability of incorrect predictions converges
exponentially to \(\approx 1\) for abstract compositional tasks.''
BABILong (Kuratov et al., 2024) reports that ``popular LLMs effectively
utilise only 10--20\% of the context and performance declines sharply
with increased reasoning complexity.'' METR (Kwa et al., 2025) reports
that the 80\%-success horizon is 4--6× shorter than the 50\%-success
horizon, suggesting a steep within-model reliability cliff. The
Fano-style upper bound of Wan et al.~(2025) explicitly theorises
super-linear information-demand growth and identifies an ``accuracy
cliff'' at capacity overflow.

The apparent tension between these two camps dissolves on careful
reading: every paper in the second group decays over a variable distinct
from raw token length \(n\). We document this in detail in Section 6.

\begin{center}\rule{0.5\linewidth}{0.5pt}\end{center}

\subsection{3. Theoretical Framework}\label{theoretical-framework}

\subsubsection{\texorpdfstring{3.0 What \(|C|\) refers to: four
levels}{3.0 What \textbar C\textbar{} refers to: four levels}}\label{what-c-refers-to-four-levels}

LLM error analysis routinely conflates four distinct objects. To keep
the rest of this section honest, we distinguish them and tag the formal
claims accordingly.

\begin{itemize}
\tightlist
\item
  \textbf{L1: failure events} --- concrete LLM errors observed in a
  benchmark. These are the raw data behind any taxonomy.
\item
  \textbf{L2: empirical taxonomy categories} --- researcher-chosen
  labels grouping L1 events. ErrorAtlas's 17 categories, MWPES's top-4
  classes, and HumanEval's \texttt{AssertionError} / \texttt{NameError}
  partition belong here. L2 is observable but depends on the taxonomer's
  resolution choices.
\item
  \textbf{L3: latent failure modes} --- the (unobserved) underlying
  clusters that L2 categories approximate. L3 is the quantity Postulate
  1 is morally about, but it cannot be measured directly; we treat L2 as
  a noisy proxy for L3.
\item
  \textbf{L4: capability axes / interventions} --- the engineering unit.
  A Python interpreter, a constrained decoder, a retrieval-augmented
  generator. L4 is coarser than L2 in practice: one capability axis
  typically targets several L2 categories at once (Section 4.2.0).
\end{itemize}

Throughout this section, \(|C|\) refers to the \textbf{L2} count --- the
number of named taxonomy categories at a given corpus scale --- because
L2 is what published taxonomies report. Postulate 1 (Section 3.2) is a
statement about how L2 grows with sampled failures. Proposition 2
(Section 3.4) is engineering math conditional on Postulate 1. The
framework is \emph{engineering}-relevant because L4 is coarser than L2
(Section 3.2 closing); it is \emph{epistemically} conditional because L2
is a proxy for the actual L3 mode catalogue, and the proxy's
faithfulness has not been independently measured.

\subsubsection{\texorpdfstring{3.1 \(\beta\)-stratification of key
tokens}{3.1 \textbackslash beta-stratification of key tokens}}\label{beta-stratification-of-key-tokens}

The two-rate model of Arbuzov et al.~(2025), reproduced as Equation (1),
distinguishes \(k\) key tokens (error rate \(e_{\text{key}}\)) from
\(n - k\) non-key tokens (error rate
\(e_{\text{non}} \ll e_{\text{key}}\)). The empirical failure-mode
atlases, however, suggest that even within the key-token class errors
are not uniformly distributed. Most key tokens correspond to
``decisions'' for which the model has stable representations and
produces correct outputs; only a fraction concentrate the actual
failures.

\textbf{Definition 1 (Hard fraction).} Partition the \(k\) key tokens of
a sequence into easy and hard subsets:

\[
k_{\text{hard}} = \beta \cdot k, \qquad k_{\text{easy}} = (1 - \beta) \cdot k, \qquad \beta \in (0,1),
\]

where hard key tokens have an elevated error rate \(e_{\text{hard}}\)
corresponding to manifold-transition decisions (in the sense of Arbuzov
et al., 2025, §3.2), and easy key tokens have
\(e_{\text{easy}} \approx e_{\text{non}}\).

The composed reliability becomes

\[
P(\text{correct}) = (1 - e_{\text{hard}})^{\beta k} \cdot (1 - e_{\text{easy}})^{(1 - \beta) k} \cdot (1 - e_{\text{non}})^{n - k}.
\]

\emph{Reading the product form.} We do not assume iid token errors. The
three rates \(e_{\text{hard}}\), \(e_{\text{easy}}\), \(e_{\text{non}}\)
summarise \emph{conditional per-decision failure hazards} after
conditioning on prior relevant decisions being correct; grouping these
hazards by stratum yields the multiplicative survival expression above.
Where within-stratum hazards vary substantially across positions, or
where error cascades dominate, the expression should be read as a
stylised reliability approximation rather than as an exact generative
model.

Because \(e_{\text{easy}} \approx e_{\text{non}}\) once context
accumulates, the failure rate is dominated by the \(\beta k\) hard
tokens, and the rest of this paper treats \(e_{\text{hard}}\) as the
load-bearing quantity.

\textbf{On the identifiability of \(\beta\).} \(\beta\) is a
\emph{latent} stratification parameter: it asks what fraction of
key-token \emph{decisions} are hard, not what fraction of \emph{observed
errors} fall into the most common named categories. Failure-mode atlases
report \(\Pr(\text{category} \mid \text{error occurred})\), which is a
different conditional from
\(\Pr(\text{hard} \mid \text{key-token decision})\). The qualitative
head-concentration reported by ErrorAtlas (17 categories sorted by
prevalence) and MWPES (dominant categories per model) and the HumanEval
two-class coverage (86.35\%) are all consistent with \(\beta\) being
non-trivial --- they show failures are concentrated when failures occur
--- but they do not pin down \(\beta\). A more identifying measurement
would require ablating model confidence at predicted decision
boundaries, which the literature has not done at scale. We therefore
carry \(\beta\) symbolically through the rest of the framework and
report all downstream conclusions parametrically.

\subsubsection{3.2 Empirical postulate: logarithmic mode
discovery}\label{empirical-postulate-logarithmic-mode-discovery}

The \(\beta k\) hard-token failures across a corpus are not arbitrary:
they cluster into a finite set of recurring failure patterns indexed
\(c = 1, \ldots, |C|\). By the stratified-manifold argument of Arbuzov
et al.~(2025, §3.2), each pattern corresponds to a specific
manifold-transition mode, so the catalogue of modes is bounded in any
fixed domain.

The natural question is how \(|C|\) grows with the number of observed
failures. Two competing predictions exist in the type--token literature:

\begin{enumerate}
\def\labelenumi{\arabic{enumi}.}
\item
  \textbf{Heaps' law (standard).}
  \(|C|(k_{\text{hard}}) \approx K \cdot k_{\text{hard}}^{b}\) with
  typical exponent \(b \in [0.4, 0.6]\) (Manning et al., 2008). This is
  the canonical type--token growth law and is asymptotically compatible
  with Zipfian rank-frequency distributions.
\item
  \textbf{Logarithmic discovery.}
  \(|C|(k_{\text{hard}}) \leq A + \sigma \cdot \ln(k_{\text{hard}})\).
\end{enumerate}

\textbf{Important: Zipfian rank-frequency does not imply logarithmic
growth.} Although it is tempting to derive logarithmic mode discovery
from a Zipfian pattern-frequency distribution, the derivation fails:
Heaps' law is the correct type--token consequence of Zipf-distributed
events, and it is power-law, not logarithmic. We therefore state
logarithmic mode discovery as an empirical postulate, defensible by
direct measurement of error taxonomies, rather than as a theorem.

\textbf{Notation: corpus vs sequence.} Two sample-size variables matter
here, and conflating them is the seam that lets reviewers in. Let
\(h(n) = \beta k(n)\) denote the number of hard-token decision
opportunities in a single sequence of length \(n\). Let \(T\) denote the
number of observed hard-failure events in a corpus used to discover or
estimate the domain's failure catalogue. These are not the same object:
\(h(n)\) controls per-sequence exposure, while \(T\) controls empirical
catalogue discovery. We further write \(C_D\) for the full reachable
catalogue of recurring failure modes inside domain patch \(D\) (defined
at the end of this section), \(C_{\text{seen},D}(T)\) for the subset
discovered after \(T\) sampled failures, and \(C_{\text{active},D}(n)\)
for the subset that a single sequence of length \(n\) can activate. We
write \(\ln T\) and \(\ln h(n)\) for clarity; small-sample bounds may be
read with \(\ln(1 + T)\) and \(\ln(1 + h(n))\), and discrete counts such
as \(\beta k\) are interpreted either as expected values or as
\(\lfloor \beta k \rfloor\) asymptotically. These reformulations do not
affect any rate claim below.

\textbf{Postulate 1 (Patch-indexed catalogue discovery).} \emph{Within a
fixed application domain \(D\), the number of named recurring failure
modes discovered after \(T\) sampled hard-failure events is bounded by}

\[
|C_{\text{seen},D}(T)| \leq \min\bigl(A_D + \sigma_D \ln T,\ |C_D|\bigr), \qquad \sigma_D > 0,\ A_D \geq 0.
\]

\emph{Here \(T\) is a corpus-level sampling variable. The cap \(|C_D|\)
is the domain-imposed ceiling: the full reachable catalogue of recurring
modes inside patch \(D\). This postulate concerns catalogue discovery,
not the number of hard decisions inside a single sequence.}

\textbf{Assumption 2 (Per-sequence active-mode exposure).} \emph{For a
single sequence of length \(n\) in domain \(D\), let
\(C_{\text{active},D}(n)\) be the set of recurring failure modes that
can be activated by its \(h(n) = \beta k(n)\) hard decisions. We assume}

\[
|C_{\text{active},D}(n)| \leq \min\bigl(A'_D + \sigma'_D \ln h(n),\ |C_D|\bigr).
\]

\emph{This is a weaker, per-sequence analogue of Postulate 1. Longer
sequences expose the model to more active failure modes only slowly,
until the domain ceiling \(|C_D|\) is reached. The primed constants
\(A'_D\), \(\sigma'_D\) are distinct from those of Postulate 1 ---
corpus discovery and per-sequence activation need not share the same
rate. Proposition 2's sequence-length scaling uses
\(C_{\text{active},D}(n)\); full domain-library budgeting uses \(C_D\).}

\textbf{Empirical calibration.} Existing LLM error taxonomies provide
endpoint anchors for small named catalogues at observed corpus scales.
ErrorAtlas (Ashury-Tahan et al., 2026), HumanEval-style code taxonomies
(Wen et al., 2024), and MWPES (Sun et al., 2025) place \(|C|\) roughly
in the \(8\)--\(20\) range across corpora ranging from approximately
\(10^4\) to \(3 \times 10^5\) observed failures, yielding
\(\sigma \in [0.87, 1.85]\) under the simple \(A=0\) logarithmic
calibration. We use \(\sigma \approx 1.85\) as a conservative planning
value, not as a fitted law. These endpoint counts are not discovery
curves: the missing test is a subsample-vs-discovered-modes measurement
within a fixed deployment patch. Full calibration details are in
Appendix C.

\textbf{Capabilities are coarser than error categories.} The category
count \(|C|\) is conservative for engineering because deployed
interventions operate at the capability level, not the label level. A
Python interpreter can remove execution-error components across
arithmetic, unit conversion, counting, list manipulation, and date
arithmetic; a constrained decoder can eliminate several structural
output failures at once. The formal coverage problem is therefore closer
to weighted set cover over failure mass than
one-intervention-per-category counting. The ranked-category model used
in Proposition 2 is a tractable proxy; the full capability harvest is in
Section 4.2.

\textbf{L2 → L4 as weighted set cover.} Formally, let each L4
intervention \(I_j\) cover a subset \(S_j \subseteq C_D\) of the L2 mode
catalogue, with per-mode residual-reduction weight \(r_{ij} \in [0, 1]\)
on mode \(i\). The deployment problem chooses a library \(\mathcal{I}\)
of size at most \(m\) to maximise covered hard-error mass:

\[
\max_{\mathcal{I} : |\mathcal{I}| \leq m}\; \sum_{i \in C_D} p_i \cdot \mathbf{1}\{i \text{ covered by some } I_j \in \mathcal{I}\}.
\]

Proposition 2 below uses the analytically tractable special case in
which the \(m\)-th-ranked intervention covers the \(m\)-th-ranked mode
(one mode per intervention, ordered by marginal covered mass). In real
deployments one L4 capability typically covers several L2 categories
simultaneously, so the ranked-category budget of Proposition 2 is a
conservative proxy for the true set-cover optimum: the empirically
achieved library size for a given residual target may be a constant
factor smaller than the formula predicts. The full set-cover treatment
of overlap, additivity, and prompt-channel interference is in Appendix
B.

\textbf{Robustness.} The qualitative claim that follows in Section 3.4
is robust to the choice between Heaps' law (power-law) and logarithmic
discovery: as long as \(|C|\) grows sublinearly in \(k_{\text{hard}}\)
\textbf{and \(k_{\text{hard}}\) is itself polylogarithmic in \(n\)} (the
\(k = \Theta(\log n)\) regime of Arbuzov et al., 2025, §3.1), the
resulting intervention-budget bound is polylogarithmic in \(n\). If
\(k_{\text{hard}}\) grows as a positive power of \(n\), the Heaps
variant inherits that power and the polylog conclusion fails along the
corresponding axis. We discuss the Heaps variant explicitly in Appendix
A.5.

\textbf{Operational patch.} A deployment patch \(D\) is more than a
topic label. It is an operational tuple fixing, over a chosen time
window, the components that determine which failure modes are reachable
at all:

\[
D = (\mathcal{X},\, \mathcal{S},\, \mathcal{U},\, \mathcal{R},\, \mathcal{E},\, \mathcal{P},\, H,\, \tau),
\]

where \(\mathcal{X}\) is the task input distribution, \(\mathcal{S}\)
the input/output schema family, \(\mathcal{U}\) the user or client
class, \(\mathcal{R}\) the retrieval corpus or knowledge source,
\(\mathcal{E}\) the evaluator or acceptance criterion, \(\mathcal{P}\)
the policy and safety constraints, \(H\) the workflow horizon
(single-turn vs multi-turn vs agentic), and \(\tau\) the time window
over which these components are fixed. A \emph{patch shift} occurs when
one or more of these components changes enough to alter the reachable
failure catalogue \(C_D\). Legal review on contracts from a fixed
jurisdiction is a different patch from legal review across mixed
jurisdictions; medical RAG over UpToDate is a different patch from
medical RAG over arbitrary web content. The patch-local claims in this
paper apply within a single, operationally fixed tuple of this form, not
to ``the domain'' in any looser sense.

\textbf{Domain patches cap the engineering problem.} A fixed deployment
patch restricts the task family, schemas, tools, workflows, and
evaluator expectations. Prior work on localised representations (Park et
al., 2024; Li \& Sarwate, 2025), cross-domain performance variation
(Wang Y. et al., 2024), and long-tail knowledge (Mallen et al., 2023;
Kandpal et al., 2023) supports the weaker claim that model behaviour is
strongly domain-dependent; it does not directly measure \(|C_D|\). We
therefore treat the patch's reachable catalogue as finite or effectively
capped as a modelling assumption motivated by domain heterogeneity, not
as a theorem derived from boundedness. Inside the patch, empirical mode
discovery is bounded by

\[
|C_{\text{observed}}(t)| \leq \min\bigl(A_D + \sigma_D \ln t,\ |C_D|\bigr),
\]

so that once enough failures have been sampled to reach \(|C_D|\),
further sampling does not enlarge the catalogue and the intervention
budget becomes a domain-constant function of \(|C_D|\) alone.
\(\sigma_D\), \(A_D\), \(\beta_D\), \(|C_D|\) are all domain-indexed.
The full evidence breakdown is in Appendix B.

Crucially, \(\sigma_D\), \(A_D\), \(\beta_D\), and \(|C_D|\) are all
domain-indexed. Deploying the same base model in three different
application domains yields three distinct mode catalogues --- not
because the model is different but because the patches are. Reliability
engineering is therefore not global manifold mastery; it is \emph{local
patch coverage}: identify the recurring transition types available in
the target patch and provision interventions against the head of that
catalogue. We return to this engineering implication in Section 5 and to
the patch-shift caveat in Section 6.3.

The discovery variable \(T\), the sequence-exposure variable \(h(n)\),
and the patch ceiling \(|C_D|\) answer different questions: \(T\)
determines how much of the catalogue we have measured, \(h(n)\)
determines how much of it a single sequence can activate, and \(|C_D|\)
determines the maximum catalogue that the domain can expose.

\textbf{Corollary 1 --- Inverse Discovery Cost.} Postulate 1 is an
\emph{upper-bound} statement: the discovered catalogue cannot grow
faster than \(A_D + \sigma_D \ln T\). Inverting the bound in the
rigorous direction yields a lower bound on the sample budget needed for
the envelope to \emph{permit} a given number of distinct modes. If
\(q > A_D\) distinct modes have been discovered after \(T\) observed
hard-failure events in patch \(D\), then under Postulate 1,

\[
T \;\geq\; \exp\!\left(\frac{q - A_D}{\sigma_D}\right).
\]

Equivalently, the logarithmic upper envelope cannot reach \(q\)
discovered modes until the sample budget is exponential in \(q\).
\emph{Tightness clause.} If the empirical discovery curve is
approximately tight against the envelope at the relevant corpus scales,
this lower bound becomes the approximate inverse cost
\(T(q) \approx \exp((q - A_D)/\sigma_D)\), and each additional
\(\Delta q\) modes raises the sample budget by a factor of approximately
\(\exp(\Delta q / \sigma_D)\) --- at the conservative calibration
\(\sigma_D \approx 1.85\), five additional discovered modes need roughly
\(15\times\) more observed hard failures, and ten additional modes need
roughly \(220\times\) more. The lower bound is unconditional under
Postulate 1; the multiplicative-cost reading is conditional on
tightness. This does not mean ordinary failures are rare; it is a
statement about \emph{newly distinguishable} tail modes, which the
envelope permits only slowly. Combined with the head-heavy coverage
model of §3.3, the economic reading is that generic open-domain training
spends increasing data on low-mass tail discovery, while patch-local
reliability directly targets the high-mass local head. Full derivation
in Appendix A.6.

\subsubsection{3.3 Coverage by a targeted intervention
library}\label{coverage-by-a-targeted-intervention-library}

Given a catalogue of \(|C|\) recurring modes, an \emph{intervention
library} of size \(m\) covers the \(m\) most prevalent modes, leaving
the remaining \(|C| - m\) uncovered. The cumulative coverage of such a
library depends on the rank-frequency distribution of the patterns
themselves.

\textbf{Empirical and approximate coverage.} Let
\(p_1 \geq p_2 \geq \cdots \geq p_{|C|}\) be the ranked empirical
hard-error masses of the local failure modes, with
\(\sum_{i=1}^{|C|} p_i = 1\). The exact cumulative coverage of the
top-\(m\) library is the empirical step function

\[
F_{\text{emp}}(m) = \sum_{i=1}^{m} p_i, \qquad m \in \{0, 1, \ldots, |C|\},
\]

with \(F_{\text{emp}}(0) = 0\), \(F_{\text{emp}}(1) = p_1\), and
\(F_{\text{emp}}(|C|) = 1\) by construction. \(F_{\text{emp}}\) handles
every edge case directly and requires no patching. For closed-form
analysis at scales where the per-mode masses are not directly measured,
we use the continuum log-head approximation

\[
F_{\log}(m; |C|) = \min\!\left(1,\ \frac{\ln m}{\ln |C|}\right), \qquad m \geq 2,\ |C| \geq 2,
\]

as a \emph{planning approximation} to \(F_{\text{emp}}\), not as the
true distribution. For our ErrorAtlas anchor (\(|C| = 17\), \(m = 5\)),
\(F_{\log}(5; 17) \approx 56.8\%\); doubling the library to \(m = 10\)
pushes it to \(F_{\log}(10; 17) \approx 81.3\%\). A Zipf-1 reference for
the same anchor gives \(H_5 / H_{17} \approx 66.4\%\), concentrating
more mass in the head than \(F_{\log}\) --- so the log form is
conservative relative to Zipf. Proposition 2 below is therefore a
closed-form planning approximation to empirical ranked coverage, not a
distribution-free theorem; the qualitative polylog conclusion is what
survives across cumulative-coverage families (Appendix A.5), and a
direct fit of \(F_{\text{emp}}\) to a measured cumulative-coverage curve
in any single patch remains an explicit falsifiability test.

A second caveat is that \(F_{\log}(m; |C|)\) tracks ranked \emph{L2
categories covered}, not \emph{L4 capability axes deployed}. A single
capability axis typically covers several L2 categories on average
(Section 4.2.0); under the set-cover model of §3.2, real-deployment
residual error at a given \(m\) is therefore often lower than
\(F_{\log}(m; |C|)\) predicts. The bound remains valid but is loose in
this regime, and empirically achieved library sizes for a given residual
threshold may be a constant factor smaller than the formula's
prediction. For \(|C| = 1\), residual error is fully covered or fully
uncovered by definition, and neither \(F_{\text{emp}}\) nor \(F_{\log}\)
is informative.

After deploying the top-\(m\) intervention library, the residual
per-hard-token error rate is

\[
e_{\text{res}}(m) = \bigl(1 - F_{\log}(m; |C|)\bigr) \cdot e_{\text{hard}}
\]

under the log-head approximation, or the corresponding
\(F_{\text{emp}}\) expression if the per-mode masses are known directly.

\subsubsection{3.4 From universal impossibility to patch-local
reliability}\label{from-universal-impossibility-to-patch-local-reliability}

The finite-catalogue claim is patch-local. It does not say that all
possible LLM failures can be covered by a universal list of fixes.
Across all possible tasks, tools, schemas, knowledge sources, workflows,
and evaluator expectations, new intervention-distinguishable failures
can keep appearing. In that setting, a finite intervention dictionary is
not a well-posed reliability target.

The positive result begins only after a deployment patch \(D\) has been
fixed. Inside a patch, the task family, schemas, tools, evaluator
expectations, and admissible workflows recur. The engineering question
then changes from ``Can one library cover all possible LLM failures?''
to ``How large must the local library be to cover enough of the
reachable failure catalogue \(C_D\)?''

\textbf{Proposition 1 (No Universal Finite Intervention Dictionary).}
\emph{Fix a residual-error tolerance \(\varepsilon\). If a domain \(D\)
contains an infinite sequence of failures where each new failure is not
covered by any finite intervention dictionary covering all earlier
failures, then the \(\varepsilon\)-resolution failure catalogue
\(C_D^{\varepsilon}\) is infinite. Consequently, no finite intervention
dictionary can guarantee residual error below \(\varepsilon\) for every
intervention-distinguishable mode in \(D\).}

\emph{Metric.} This is a worst-case, mode-wise guarantee, not a claim
about expected residual error under a fixed probability distribution.
Under a distributional metric, an infinite uncovered tail may still have
arbitrarily small total mass; the proposition rules out only the
worst-case-mode-coverage reading of universal reliability, which is the
reading the engineering literature implicitly assumes when it asks for
``a fixed list of interventions that covers LLM use.''

\emph{Proof sketch.} Each new failure is intervention-distinguishable
from the failures before it; the sequence therefore generates infinitely
many distinct intervention modes, and any finite dictionary must miss
some later mode. The full proof, including the precise definition of
mode-level coverage, is in Appendix A.1.

\textbf{Engineering meaning.} Universal reliability is not a
finite-library problem. A fixed list of interventions cannot cover
open-ended LLM use in general. The rest of the paper is therefore about
\emph{patch-local} reliability, not universal reliability. Proposition 1
is the inoculation against a misread: we are not claiming a fixed list
of ≈50 named patterns covers LLM use in general.

The positive result composes \(\beta\)-stratification, the active-mode
bound (Assumption 2), and the log-coverage form into a local budget
statement. The result is conditional engineering math: it depends on the
log-coverage approximation, the active-mode exposure assumption, and the
patch hypothesis that \(C_D\) is finite or effectively capped. Appendix
A.2 states these conditions explicitly before the derivation.

\textbf{Proposition 2 (Patch-Local Sufficient Intervention Budget).}
\emph{Fix a deployment patch \(D\). Let
\(\varepsilon \in (0, e_{\text{hard}})\) be the target per-hard-decision
residual error rate. Under the log-head approximation \(F_{\log}\) of
§3.3, a library covering the dominant local modes is sufficient for the
target once}

\[
m \;\geq\; \left\lceil |C_{\text{eff}}|^{1 - \varepsilon / e_{\text{hard}}} \right\rceil,
\]

\emph{where \(C_{\text{eff}}\) is either the active catalogue
\(C_{\text{active},D}(n)\) touched by a single sequence, or the full
reachable catalogue \(C_D\) of the deployment patch. If
\(C_{\text{active},D}(n)\) grows logarithmically with the number of hard
decisions and \(k(n) = \Theta(\log n)\), then the pre-cap sufficient
budget grows doubly-logarithmically in sequence length. Once the patch
catalogue saturates, the sufficient budget becomes domain-constant in
\(|C_D|\). This is a sufficient, model-implied budget under the stated
approximation, not a proof of the true minimal intervention library.}

\emph{Proof sketch.} The library removes cumulative hard-error mass
\(F(m; |C|)\); the residual hard-token error rate is therefore
\((1 - F)\,e_{\text{hard}}\). Requiring this residual to be at most
\(\varepsilon\) rearranges, after substituting the log-coverage form, to
the stated bound. The full step-by-step derivation, including the
pre-cap and cap regimes, is in Appendix A.2.

\textbf{Engineering meaning.} Once the patch is fixed, the problem
changes. The relevant engineering question is no longer whether
arbitrary future failures exist, but how quickly the reachable local
catalogue is discovered and how much hard-error mass is removed by the
top interventions. The intervention budget for per-hard-token
reliability is small and slowly growing, domain-constant in the cap
regime, with the exact size determined by the local discovery and
rank-coverage curves rather than by a universal prior.

\textbf{Sequence-level caveat.} Proposition 2 bounds residual error per
hard decision. A one-shot sequence-level target is strictly stricter
because many hard decisions occur in one output: as \(k\) grows the
allowable residual per hard decision shrinks, and the required library
approaches full-catalogue coverage. In some regimes the non-hard-token
error mass alone already exceeds the sequence-level budget, so
hard-token interventions cannot meet the SLA by themselves. The full
three-regime analysis (and the per-hard-token tolerance
\(\tau_{\text{seq}}\) that converts a sequence target into the
Proposition 2 bound) is in Appendix A.3.

A common engineering rule of thumb --- ``cover the local head with on
the order of tens of interventions'' --- is best understood as a
per-hard-decision planning prior, not a derived constant from the formal
model. The per-hard-decision residual-error reduction implied by
\(F_{\log}\) at typical \(|C|\) and head-heavy mass is large; the
sequence-level analogue requires correspondingly more interventions and
tighter coverage of the tail.

We call these statements propositions rather than theorems because their
force is conditional: they formalise the consequences of the paper's
modelling assumptions rather than deriving a universal law of LLM
reliability. The doubly-logarithmic rate is the optimistic special case
under logarithmic active-mode exposure; Appendix A.5 gives the Heaps and
saturation variants. Appendix A.4 expands on the conditional reading.

\begin{center}\rule{0.5\linewidth}{0.5pt}\end{center}

\subsection{4. Empirical Evidence}\label{empirical-evidence}

We now collect quantitative evidence for the three load-bearing claims:
errors cluster into a small recurring set (Section 4.1); each cluster is
independently addressable by one targeted intervention (Section 4.2);
and reliability decays sublinearly, not exponentially, with output
length (Section 4.3).

\subsubsection{4.1 Claim A: Failure-mode
clustering}\label{claim-a-failure-mode-clustering}

\textbf{Cross-domain concentration.} The strongest single result is
ErrorAtlas (Ashury-Tahan et al., 2026). Across 83 models and 35
datasets, with approximately 7,000 sampled failures, the taxonomy
resolves to 17 named categories sorted by prevalence (Table 1 of the
source). The published category set is fixed in advance, and the
empirical prevalence distribution is heavily long-tailed: a handful of
categories --- including ``Missing Required Element'', ``Logical
Reasoning Error'', and ``Computation Error'' --- dominate the failure
mass, while the remaining categories account for progressively smaller
residuals. A measured cross-domain cumulative top-\(m\) coverage curve
has not yet been published by the source; this is the explicit empirical
test we flag for the log-coverage form of §3.3.

\textbf{Domain-specific Paretos.} The cross-domain picture is reproduced
at finer grain within each domain:

\begin{itemize}
\tightlist
\item
  \textbf{Math.} MWPES-300K (Sun et al., 2025) categorises 304,865
  errors from 15 LLMs across four math word-problem datasets. The
  published per-model error breakdowns place per-category prevalence in
  single-digit-to-mid-teens-percent ranges (e.g., ``Misunderstanding of
  problem requirements'' at 4.83--14.60\% across models;
  ``Miscalculation during algebraic manipulation'' at 5.03--13.18\%); a
  small set of dominant categories accounts for the bulk of mass on each
  model, and ``dataset characteristics significantly shape error
  patterns'' (Sun et al., 2025). Whether the \emph{identical} top-\(k\)
  category set recurs across all 15 models, and the exact cross-model
  cumulative coverage of that top-\(k\) set, are claims we leave to the
  source paper's tables rather than re-state.
\item
  \textbf{Code.} Wen et al.~(2024) categorise 12,837 errors from 14 LLMs
  on HumanEval. \texttt{AssertionError} and \texttt{NameError} alone
  capture 86.35\% of failures. The dominant-category hierarchy
  replicates across 23 models on HumanEval Pro and MBPP Pro.
\item
  \textbf{Multi-hop QA.} The Weakest Link Law (Zhang et al., 2026) shows
  multi-hop performance on MuSiQue and NeoQA collapsing to the
  least-visible-evidence bucket --- a single dominant failure mode
  (missing evidence retrieval) accounting for most errors across
  question types.
\item
  \textbf{Agentic / tool use.} MAST (Cemri et al., 2025) and
  \(\tau\)-bench (Yao et al., 2024) document recurring failure modes in
  tool-augmented agents (tool selection, parameter disambiguation,
  intermediate state loss). The mode count is small (under twenty in
  each study) and the cumulative coverage of the top modes is high.
\item
  \textbf{Retrieval-augmented generation.} Acurai (Wood \& Forbes, 2024)
  on RAGTruth and follow-on clinical/legal RAG work all report that
  hallucinations cluster into a handful of distinguishable modes
  (citation hallucination, context-stitching errors, retrieval-context
  mismatch), each independently addressable.
\end{itemize}

\textbf{Cross-model invariance.} The same top categories appear across
very different models. ErrorAtlas reports a stable category hierarchy
across 83 models. RFMDataset (Guo et al., 2025) explicitly tests
invariance and reports ``strikingly similar failure mode distributions''
across 10 advanced reasoning models (three open-source and seven
proprietary). EDIT (Dai et al., 2025) finds an explicit
\(\approx 4.7\%\) key-step fraction across reasoning chains and across
models, with stable category labels.

\textbf{Theoretical support.} Schaeffer et al.~(2025) prove that the
observed aggregate power-law scaling of LLM evaluations requires the
per-problem success-rate distribution to be heavy-tailed near \(p = 0\).
This is not a direct theorem about cluster-count growth (which is the
object of our Postulate 1), but it is a closely related structural
result: the failure side of LLM behaviour is intrinsically heavy-tailed.

The aggregate picture is what Postulate 1 of Section 3.2 is designed to
capture: LLM errors are not a long uniform tail but a short, repeating
list. The ErrorAtlas \(|C| = 17\) with a long-tailed, head-concentrated
prevalence ordering is the cleanest summary; per-domain Paretos show the
same shape inside each domain; the cross-model invariance shows the
catalogue is stable.

\subsubsection{4.2 Claim B: Cluster-selective capability
interventions}\label{claim-b-cluster-selective-capability-interventions}

The intervention literature provides at least one targeted
countermeasure for each named cluster of Section 4.1, almost always with
a published quantitative result. Two organising granularities exist in
this evidence, and we report both. At the \emph{capability} level (six
axes; Section 4.2.0) the cluster-orthogonality property is most clearly
visible --- capability provisioning leaves a class-shift signature that
says exactly which class was erased and which classes the residuals
belong to. At the \emph{error-class} level (twelve named clusters;
Sections 4.2.1--4.2.2) the evidence aligns with the taxonomies of
Section 4.1 but is finer-grained than the underlying capability
mechanisms. The two views are complementary: the axes are the
engineering unit, the clusters are the bookkeeping unit.

\paragraph{4.2.0 Six capability axes and three structural
patterns}\label{six-capability-axes-and-three-structural-patterns}

A dedicated harvest of capability-elimination evidence (parallel to the
cluster-based harvest below) yields \textbf{28 quantitatively-anchored
citations} across \textbf{six independent capability axes}: Arithmetic
(Python/symbolic execution), Code Execution (REPL/sandbox feedback),
Format/Structure (constrained decoding, FSMs, grammar engines),
Perception/Grounding (visual grounding models for GUI, charts, tables),
Knowledge/RAG (inference-time dense retrieval and citation grounding),
and Verification (formal proof checkers, learned verifiers, classifier
rerouting, process supervision). Each axis is independently confirmed by
between three and nine citations. The 28 results exhibit three
structural patterns --- by-construction,
strong-empirical-with-class-shift, and moderate-with-residual-shift ---
that strengthen and refine the cluster-orthogonality claim.

We stratify these into three tiers by the \emph{kind} of evidence each
provides, taking care not to inflate empirical reductions into
mathematical guarantees.

\textbf{Pattern A --- Hard guarantees (by-construction).} Seven
citations achieve residual error rate equal to zero \emph{by
construction}, restricted strictly to structural / verifiable error
classes: constrained decoders set \(P(\text{invalid token}) = 0\) at
every generation step (DINGO, ToolDec, XGrammar, XGrammar-2, OpenAI
Structured Outputs); a static syntax check rejects any program with a
\texttt{SyntaxError} before execution (LlmFix's \texttt{SyntaxError}
repair, 5.76\% → 0.01\%); proof kernels reject any output that fails
type-checking (Lean / DeepSeek-Prover-V2). The class of ``outputs that
violate the grammar'' is mathematically empty under these mechanisms. We
restrict Pattern A to this narrow sense: format, syntax, schema,
formal-proof validity. (LlmFix's \texttt{NameError} row, by contrast, is
a Pattern-C statistical reduction --- 22.7\% → 2.3\%, not zero --- and
is tabulated as such.)

\textbf{Pattern B --- Strong empirical reductions with class-shift
signature.} Roughly fourteen citations report empirical reductions of
80--100\% in a named error class, with the post-intervention failure log
dominated by structurally different residual classes. The cleanest
example is Program-of-Thoughts on GSM8K (Chen et al., 2022): manual
taxonomy of 50 PoT failures shows calculation errors drop from 30\% of
CoT failures to \textbf{0\%}, while residuals are 62\% reasoning + 36\%
misunderstanding. In the OpenMedCalc arm of Goodell et al.~(2025),
\emph{``only interpretation errors were identified''} --- the
calculation-error class is empirically absent from the post-tool failure
log. Acurai (2024) reports 100\% hallucination elimination on RAGTruth
(95\% CI 91--100\%) for GPT-4 and GPT-3.5 --- a strong empirical result
on that benchmark, not a by-construction property of RAG in general. The
class-shift signature is the structural observation: post-intervention
residuals belong to different classes, consistent with cluster
orthogonality, though formally the framework only requires this in
expectation rather than identically.

\textbf{Pattern C --- Moderate reductions (60--80\%) with residuals
shifting outside the target class.} Several citations: Legal RAG
(Dantart, 2026) reduces fabricated citations from \textgreater30\% to
\textless0.2\%; enabling web search on GPT-5 reduces SimpleQA factual
errors from 47\% (offline) to 9.6\% (with web access enabled) --- an
inter-condition comparison, not a within-condition before/after (OpenAI,
2025); CRITIC toxicity (Gou et al., 2024) reduces toxic generation by
79.2\%. Residuals shift toward reasoning / planning errors rather than
inflating the targeted class. This is the regime where the framework's
``selectivity'' claim is supported but does not rise to elimination.

The tier boundaries matter for what the paper claims. Pattern A admits
the strongest rhetoric: ``by construction,'' ``mathematically empty.''
Pattern B admits ``strong empirical reduction'' and ``class-shift
signature'' --- but not ``by construction''; Acurai's 100\% on RAGTruth
is a benchmark-conditional empirical fact, not a mathematical guarantee.
Pattern C supports ``selective reduction'' without claiming class
elimination. Keeping these tiers crisp is what saves the paper from
drifting into product-brochure language about capability provisioning.

The 28 citations, organised by axis and pattern, appear below.

{\def\LTcaptype{none} % do not increment counter
\begin{longtable}[]{@{}
  >{\raggedright\arraybackslash}p{(\linewidth - 8\tabcolsep) * \real{0.2000}}
  >{\raggedright\arraybackslash}p{(\linewidth - 8\tabcolsep) * \real{0.2000}}
  >{\raggedright\arraybackslash}p{(\linewidth - 8\tabcolsep) * \real{0.2000}}
  >{\raggedright\arraybackslash}p{(\linewidth - 8\tabcolsep) * \real{0.2000}}
  >{\raggedright\arraybackslash}p{(\linewidth - 8\tabcolsep) * \real{0.2000}}@{}}
\toprule\noalign{}
\begin{minipage}[b]{\linewidth}\raggedright
Axis
\end{minipage} & \begin{minipage}[b]{\linewidth}\raggedright
Citation
\end{minipage} & \begin{minipage}[b]{\linewidth}\raggedright
Targeted class
\end{minipage} & \begin{minipage}[b]{\linewidth}\raggedright
Pre → Post
\end{minipage} & \begin{minipage}[b]{\linewidth}\raggedright
Pattern
\end{minipage} \\
\midrule\noalign{}
\endhead
\bottomrule\noalign{}
\endlastfoot
Verification & DeepSeek-Prover-V2 (Ren et al., 2025) & Invalid Lean
proof & any → 0\% by construction & A \\
Format & OpenAI Structured Outputs (2024) & JSON schema violation &
\textgreater60\% → 0\% by construction & A \\
Format & ToolDec (Zhang et al., 2023) & Tool-call syntax error &
21--100\% → 0\% & A \\
Format & DINGO (Suresh et al., 2025) & JSON parse failure & 13--82\% →
0\% & A \\
Format & XGrammar (Dong et al., 2024) & Multi-format errors & 20--38\% →
0\% & A \\
Format & XGrammar-2 (Li et al., 2026) & Malformed tool calls & 33--78\%
→ 0\% & A \\
Arithmetic & Program of Thoughts (Chen et al., 2022) & Calculation
errors on GSM8K & 30\% of failures → 0\% & B \\
Arithmetic & OpenMedCalc (Goodell et al., 2025) & Clinical arithmetic &
``only interpretation errors remained'' & B \\
Knowledge/RAG & Radiology RAG (Wada et al., 2025) & RAG hallucinations &
8\% → 0\% (p=0.012) & B \\
Knowledge/RAG & Acurai (2024) & Context-conflict hallucinations & 100\%
→ 0\% on conflict subset & B \\
Knowledge/RAG & Rigorous Archivist (Dantart, 2026) & Fabricated legal
citations & \textgreater30\% → \textless0.2\% & C \\
Code Exec & LlmFix (Wen et al., 2024) & SyntaxError (HumanEval) & 5.76\%
→ 0.01\% & A \\
Code Exec & AlphaCode (Li et al., 2022) & False-positive submissions &
62\% → 4\% & B \\
Code Exec & MGDebugger (Shi et al., 2024) & 5 code-bug classes & 100\%
repair (5/6 classes) & B \\
Knowledge/RAG & ALCE (Gao et al., 2023) & Citation grounding (ELI5) &
best models leave \textasciitilde50\% of statements without complete
citation support & C \\
Knowledge/RAG & GPT-5 SimpleQA (OpenAI, 2025) & Factual errors & 47\%
offline → 9.6\% with web access enabled & C \\
Knowledge/RAG & Almanac (NEJM AI 2024) & Clinical citation errors & 44\%
→ 9\% & C \\
Arithmetic & MedRaC (Wang B. et al., 2025) & Arithmetic in medical
reasoning & 426 → 74 errors & B \\
Code Exec & LlmFix NameError (Wen et al., 2024) & NameError repair &
22.7\% → 2.3\% & C \\
Perception & UGround (Gou et al., 2025) & GUI grounding errors & 16.2\%
→ 73.3\% accuracy & B \\
Perception & Jedi / OSWorld (Xie et al., 2025) & GUI task failures & 5\%
→ 27\% SR (5.4×) & B \\
Perception & SeeClick (Cheng et al., 2024) & Element mis-location &
5.2\% → 53.4\% (10.3×) & B \\
Perception & DePlot (Liu et al., 2023) & Chart-reading errors & 38.2\% →
67.6\% (+29.4 pp) & B \\
Verification & CRITIC toxicity (Gou et al., 2024) & Toxic generation &
79.2\% reduction & C \\
Verification & Math-Shepherd (Wang P. et al., 2023) & Reasoning step
errors on MATH (Mistral-7B) & 28.6\% → 33.0\% (PPO); 43.5\% with
verification reranking & B \\
Knowledge/RAG & Self-RAG (Asai et al., 2024) & Unsupported facts &
55.5\% → 22\% error & B \\
Perception & OmniParser (Lu et al., 2024) & Icon / element errors &
70.5\% → 93.8\% (79\% error reduction) & B \\
Knowledge/RAG & PopQA (Mallen et al., 2023) & Long-tail entity errors &
\textasciitilde80\% → \textasciitilde50\% failure (\textasciitilde75\%
on tail) & B \\
\end{longtable}
}

\textbf{Capabilities are coarser than error classes.} Several entries in
the harvest reveal ``two-for-one'' reductions where one capability
addresses multiple named clusters. A single Python interpreter removes
the execution-error component of arithmetic (cluster A), unit conversion
(B), simple counting (C), list manipulation, and date arithmetic ---
five named classes under one provisioning, with residual failures
shifting to problem representation and code synthesis. Constrained
decoding eliminates by construction both format violations (D) and the
structural component of ``missing required elements'' (one of the most
prevalent ErrorAtlas categories, our Category I), because schemas with
\texttt{required:\ {[}...{]}} fields cannot omit those fields. Code
execution feedback strongly reduces SyntaxError, NameError, and most
TypeError (three subclasses of cluster E) together. RAG strongly reduces
factual hallucinations (J), fabricated citations, and outdated
information jointly when the relevant evidence is retrieved and
grounded, while leaving residual failures in retrieval coverage, context
selection, and reasoning over retrieved material. The practical
capability library required to address a given fraction of failures is
therefore smaller than the count of named error categories --- the
empirical phenomenon noted in Section 3.2 as a conservative-bound
tightening of Postulate 1.

A negative-control result confirms the framework's selectivity:
DebugBench (Tian et al., 2024) finds that execution feedback works well
for syntax/reference errors but is \emph{``unhelpful for logic
errors\ldots{} may even cause disruptions.''} Capability provisioning
addresses what it can detect or constrain, and the prediction of where
it fails (semantic / open-ended classes) is correct. We return to this
point in Section 6.

\paragraph{Summary of the twelve
categories}\label{summary-of-the-twelve-categories}

{\def\LTcaptype{none} % do not increment counter
\begin{longtable}[]{@{}
  >{\raggedright\arraybackslash}p{(\linewidth - 12\tabcolsep) * \real{0.1429}}
  >{\raggedright\arraybackslash}p{(\linewidth - 12\tabcolsep) * \real{0.1429}}
  >{\raggedright\arraybackslash}p{(\linewidth - 12\tabcolsep) * \real{0.1429}}
  >{\raggedright\arraybackslash}p{(\linewidth - 12\tabcolsep) * \real{0.1429}}
  >{\raggedright\arraybackslash}p{(\linewidth - 12\tabcolsep) * \real{0.1429}}
  >{\raggedright\arraybackslash}p{(\linewidth - 12\tabcolsep) * \real{0.1429}}
  >{\raggedright\arraybackslash}p{(\linewidth - 12\tabcolsep) * \real{0.1429}}@{}}
\toprule\noalign{}
\begin{minipage}[b]{\linewidth}\raggedright
Cluster
\end{minipage} & \begin{minipage}[b]{\linewidth}\raggedright
Covering axis
\end{minipage} & \begin{minipage}[b]{\linewidth}\raggedright
Intervention
\end{minipage} & \begin{minipage}[b]{\linewidth}\raggedright
Before → After
\end{minipage} & \begin{minipage}[b]{\linewidth}\raggedright
\(\Delta\)
\end{minipage} & \begin{minipage}[b]{\linewidth}\raggedright
Strength
\end{minipage} & \begin{minipage}[b]{\linewidth}\raggedright
Citation
\end{minipage} \\
\midrule\noalign{}
\endhead
\bottomrule\noalign{}
\endlastfoot
A. Arithmetic / computation & Arithmetic & Python interpreter (PAL) &
20.1\% → 61.5\% & +41.4 pp & Strong & Gao et al., 2022 \\
B. Unit conversion & Arithmetic & Folded into A (Python handles units
natively) & --- & --- & (Folded) & --- \\
C. Counting / enumeration & Arithmetic & Explicit counter (research gap)
& --- & --- & Gap & Mechanistic interp. work \\
D. Format / schema violation & Format/Structure & Constrained decoding
(DINGO, SLOT) & up to +68 pp (math+JSON benchmarks); 99.5\% schema acc.
& +68 pp & Strong & Suresh et al., 2025; Wang D. Y.-B. et al., 2025 \\
E. Code logic / type / null & Code Exec & REPL feedback (Reflexion,
AgentCoder) & 80\% → 96.3\% pass@1 & +16.3 pp & Strong & Shinn et al.,
2023 \\
F. Multi-hop decomposition drift & Knowledge/RAG & Entity-grounded
rewriter (ChainRAG) & consistent gains across 3 benchmarks &
\$\sim\$5--10 pp & Medium & Zhu et al., 2025 \\
G. Reasoning step errors & Verification & Process reward model
(Math-Shepherd) & Mistral-7B MATH: 28.6\% → 33.0\% (PPO); → 43.5\% with
verification reranking & +4.4 / +14.9 pp & Strong & Wang P. et al.,
2023 \\
H. Specification misinterpretation & (semantic residual) & Clarification
loops (AmbigNLG) & ROUGE-L +15; compliance 86\% → 94\% & +8 pp & Medium
& Niwa \& Iso, 2024 \\
I. Missing required elements & Format/Structure & Reframed as Category D
for structural, G/H for semantic & (subsumed) & --- & Reframed & This
work \\
J. Hallucination / factual errors & Knowledge/RAG & RAG query
reformatting (Acurai) & non-zero → 0\% & \(-100\%\) & Strong & Wood \&
Forbes, 2024 \\
K. Inappropriate refusal & Verification & Preference optimisation
(POROver) & 57.6\% → 82.1\% & +24.5 pp & Strong & Karaman et al.,
2024 \\
L. Tool / API usage errors & Format/Structure + Verification &
Structured uncertainty (SAGE-Agent) & 36.5\% → 65.2\% & +28.7 pp &
Strong & Suri et al., 2025 \\
\end{longtable}
}

Eight of twelve categories have a strong citation with double-digit
percentage-point improvement; three (B, C, I) lack a clean
single-cluster ablation in the literature and are addressed below; one
(F) has consistent medium-strength evidence.

\paragraph{Detailed evidence on the strongest
categories}\label{detailed-evidence-on-the-strongest-categories}

The arithmetic category is the cleanest example of cluster
orthogonality. Delegating execution to a Python interpreter, PAL (Gao et
al., 2022) lifts GSM-Hard accuracy from 20.1\% to 61.5\% --- a
41.4-point gain on the benchmark configuration where large numbers
stress LLM tokenisation. What matters more than the gain is the
composition of the residual: PAL's remaining failures are
problem-comprehension errors, mis-translations from word problem to
code. Arithmetic execution errors have left the failure log. Toolformer
(Schick et al., 2023) self-teaches the same shift and lands a 97.9\%
calculator-call rate on arithmetic queries; and on clinical calculations
Goodell et al.~(2025) report a 5--13× error reduction (LLaMA: 16\% →
88\%; GPT-4: 4.8\% → 64\%) once the appropriate domain tool is in the
loop.

Format violations admit something stronger still: a mathematical
guarantee. DINGO (Suresh et al., 2025) reports up to +68 percentage
points on combined symbolic-math and JSON-generation benchmarks. SLOT
(Wang D. Y.-B. et al., 2025) reaches 99.5\% schema accuracy with 94\%
content similarity --- the cluster is essentially gone. A systematic
evaluation of four production frameworks (Guidance, Outlines, Llamacpp,
XGrammar) by JSONSchemaBench (Geng et al., 2025) reports coverage up to
0.96 on real-world schemas. The mechanism is uniform across the family:
zero out the logits of grammar-violating tokens at each decoding step
and the violations cannot occur.

The hallucination result is the most striking single number in the
harvest. Acurai (Wood \& Forbes, 2024) --- query reformatting applied to
RAG --- reports 100\% hallucination elimination on RAGTruth for both
GPT-4 and GPT-3.5, with a 95\% CI of 91--100\%. This is a
benchmark-specific empirical result, not a general by-construction
guarantee: Acurai does not prove that hallucinations cannot occur, only
that none were observed on RAGTruth's test set under their specific
query-reformatting setup. Domain-specific variants (legal, clinical)
report 60--90\% reduction with magnitudes depending on the baseline. The
mechanism is austere but not mathematical: hallucinations require the
model to generate factual claims absent from context, and
citation-grounded retrieval tightens --- without eliminating in general
--- that failure path.

Code generation shows a clean staircase. Reflexion (Shinn et al., 2023)
lifts HumanEval pass@1 from 80\% to 91\% by reflecting on execution
feedback; AgentCoder (Huang et al., 2023) extends the loop to
multi-agent test-driven generation and reaches 96.3\% --- approaching
ceiling. At 96.3\% the residual is dominated by spec-misinterpretation
(Category H). The execution-feedback capability has erased what it can
erase.

\paragraph{The Category-I reframing}\label{the-category-i-reframing}

Of the three apparent gaps (B, C, I), Category I --- ``missing required
elements'' --- is the most significant because it is one of the most
prevalent named categories in ErrorAtlas. Originally we expected to find
a dedicated ``checklist verifier'' intervention in the literature; we
did not. The cleanest resolution is mechanistic equivalence with
Category D. A ``missing required element'' error is structurally a
schema-violation where the schema's \texttt{required:\ {[}...{]}} field
list is not satisfied. Constrained decoders (DINGO, SLOT, Outlines,
Guidance, XGrammar) enforce schema constraints during decoding; with a
properly specified schema, they mathematically \emph{cannot} produce
output that omits a required field. SLOT's 99.5\% schema accuracy
directly includes \texttt{required} enforcement.

We therefore split Category I into two mechanistically distinct
sub-clusters: (I.a) structural missing elements --- omitted JSON field,
missing function argument, missing required citation --- which are
absorbed by Category D's constrained-decoding interventions; and (I.b)
semantic missing elements --- incomplete reasoning step, missing edge
case in code, omitted assumption in proof --- which fall under Category
G (process supervision) or Category H (clarification loops). No new
intervention category is required, and two of the highest-volume
clusters (D format + I.a structural missing, jointly \(\approx 25\%\) of
ErrorAtlas) share a single mechanism --- which strengthens rather than
weakens the cluster-orthogonality argument.

Categories B (unit conversion) and C (counting) remain technical gaps.
Unit conversion is naturally folded into Category A: any Python
interpreter handles dimensional analysis, and PAL/PoT/ToRA results
already encompass it. Counting is the smallest residual gap; a focused
literature pass or a one-day inference-only ablation would close it.

\paragraph{Additivity and its limits}\label{additivity-and-its-limits}

Six Sigma Agent results (Patel et al., 2026) demonstrate that stacking
interventions targeting orthogonal failure modes can produce large
compound reliability gains. In their parallel-consensus framework,
atomic task decomposition combined with five-way majority voting yields
a \(14{,}700\times\) improvement over a single-pass baseline on the
authors' internal benchmark --- we read this as evidence for compound
gains from decomposition plus consensus aggregation under their specific
parallel-voting regime, not as a direct demonstration that heterogeneous
cluster-targeted interventions compose additively in the absence of
voting. AgentSquare (Shang et al., 2024) similarly shows supra-additive
gains when modules target distinct failure modes. These are encouraging
but benchmark- and aggregation-specific.

There is a caveat. Le (2026) reports that schema-level and prompt-level
instructions interact \emph{non-additively} for some model families:
stacking constrained decoding (Category D) with prompt-level
specification clarifications (Category H) produces less than the sum of
individual gains, sometimes substantially so. The interaction is
concentrated in cases where both interventions share a common prompt
channel. We treat this as evidence not against additivity in general but
for a refinement: interventions that operate on orthogonal processing
layers (decoding constraint vs.~retrieval vs.~training-signal
vs.~inference-time tool call) compose near-additively, while those that
share a channel may interfere. This refinement is itself testable and we
leave it as future work.

\subsubsection{4.3 Claim C: Sublinear length
scaling}\label{claim-c-sublinear-length-scaling}

We now compile direct evidence that LLM reliability decays sublinearly
--- typically logarithmically, sigmoidally, or in a plateau-then-cliff
pattern --- rather than exponentially with output or context length.
This is the headline conclusion of Arbuzov et al.~(2025); we extend it
here with newly available benchmark results and reconcile it with the
steeper-decay reports re-examined in Section 6.

\textbf{Log-linear decay on multi-document reasoning.} Loong (Wang M. et
al., 2024) evaluates 4 task types on 4 context-length sets (10K--50K,
50K--100K, 100K--200K, \textgreater200K). GPT-4o on Chain-of-Reasoning
declines 81.6\% → 62.1\% → 45.1\% → 32.9\% across the four sets ---
consistent with log-linear decay across the measured bins. Exponential
per-token decay would predict near-zero accuracy by 200K and is
decisively ruled out by the 32.9\% endpoint. We do not commit to a
specific log-linear fit constant because the bin-representative
convention (lower bound, midpoint, mean) materially shifts the slope;
the qualitative claim --- mild, much-slower-than-exponential decay over
\(\log L\) --- is robust to that choice.

\textbf{Sigmoid decay on synthetic math.} Exponentially more inference
compute, on GSM-\(\infty\), yields only \emph{linear} AUC gains
(\(R^2 > 0.99\) for the linear fit) --- under independent per-step
Bernoulli errors, more trials should give exponential improvement. The
same benchmark fits a sigmoid (saturating) decay with \(R^2 > 0.98\) on
the Medium subtask, and DeepSeek-R1 maintains 10\%+ accuracy at 130
reasoning operations where \((1 - \varepsilon)^{130}\) with any
reasonable \(\varepsilon\) would predict near-zero (Zhou et al., 2025).
Errors concentrate at semantic decision points; the data refuse to
behave like uniformly-distributed token-level noise.

\textbf{Compositionality gap is bounded, not growing.} The gap between
two-hop accuracy and the product of one-hop accuracies should, under
independent per-step compounding, grow with chain length; under
clustered errors, it should not. Press et al.~(2023) find a
\(\sim 40\%\) gap that is approximately \emph{constant} across the GPT-3
family (sub-billion to 175B parameters). Independent compounding is
directly ruled out.

\textbf{Self-consistency is consistent with --- but does not prove ---
clustered structure.} Wang X. et al.~(2023) report +17.9 pp on GSM8K
from majority voting over 40 sampled chains. A common reading interprets
the magnitude as incompatible with independent sampled paths, but this
is not quite right: Condorcet aggregation of iid Bernoulli(\(p\)) chains
with \(p > 0.5\) can move accuracy substantially as the sample count
grows, so gain magnitude alone is not diagnostic. The
clustered-attractor interpretation remains \emph{consistent} with the
observed gain --- sampled chains converge on a small number of recurring
reasoning paths, exactly the structural prediction of Section 3 --- but
is not established by it. We therefore retain self-consistency as
suggestive evidence pointing in the same direction as Section 4.1's
concentration result, not as an independent proof of cluster structure.

\textbf{Bounded effective decision count via K/V compression.} Pang et
al.~(2024)'s Anchor LLMs achieve \(\approx 99\%\) reduction in cached
K/V state per layer with only minor accuracy compromise, plus a 3.5×
inference speedup. The mechanism operates on internal K/V vectors, not
on the input token stream --- so the observation does not directly imply
input-token redundancy. What it does imply is that the \emph{effective
number of distinct attention-key vectors} a long-context model needs to
retain to produce coherent output is far smaller than \(n\) --- exactly
the architectural form of the sparse-decision claim. LongLLMLingua
(Jiang et al., 2024) and Recurrent Context Compression (Huang et al.,
2024) report similar prompt- and context-compression bounds along the
same theme.

\textbf{Staircase / plateau-and-cliff patterns.} Think-Prune-Train
(Costello et al., 2025) reports that successive self-training rounds
raise Pass@1 (Gemma-2B: 41.9\% → 57.6\%; LLaMA-70B: 78\% → 91.5\%) while
Pass@20 \emph{plateaus} after the first round. This is the
manifold-transition signature: the model is finding the correct cluster
more reliably (Pass@1 rises) but is not finding new correct solutions
(Pass@20 unchanged). RULER (Hsieh et al., 2024) reports threshold
behaviour: simple retrieval holds to 32K; complex reasoning collapses
earlier. Neither pattern is consistent with smooth
\((1 - \varepsilon)^n\) decay.

\textbf{METR's logistic-in-log-length.} METR (Kwa et al., 2025) fits
per-model success as a logistic function of log(human-time):
\(S(t) = \sigma(\beta \cdot (\log t - \log h))\). Logistic-in-log is
mathematically sublinear in length itself. The famous METR exponential
--- capability doubling every seven months --- is an \emph{inter-model}
claim about how the horizon \(h\) grows over generations, not an
intra-model claim about within-model decay shape. We address the common
misreading in Section 6.

\textbf{Counter-evidence.} Several prominent papers report what reads as
rapid or ``exponential'' decay --- Dziri et al.~(2023), Kuratov et
al.~(2024), METR (2025), Wan et al.~(2025). Read carefully, each decays
over a variable other than raw token length: compositional graph size,
number of supporting facts, log-time horizon, capacity threshold.
Section 6 walks through each one.

\begin{center}\rule{0.5\linewidth}{0.5pt}\end{center}

\subsection{5. Practical Implications}\label{practical-implications}

\textbf{Reliability engineering is local patch coverage.} Within a fixed
domain patch, Proposition 2's polylog bound makes reliability a
small-catalogue engineering problem rather than an asymptotic scaling
problem. A team building a reliable LLM system in a measured domain
should initially budget for a library on the order of tens of
interventions, then refine that estimate from the local mode-discovery
curve \(C_{\text{seen},D}(T)\) and the empirical rank-frequency
distribution. As a rule of thumb consistent with the head-mass figures
in ErrorAtlas, HumanEval, and MWPES, \(\approx 50\) named interventions
cover the bulk of the per-hard-token failure distribution in many
measured domains. This is a planning prior calibrated to current
taxonomies, not a universal constant; the local mode-discovery curve
sets the actual library size for any given deployment. The same base
model in three different patches (cardiology RAG, legal contract
drafting, code review) yields three different intervention libraries:
\(C_D\), \(A_D\), and \(\sigma_D\) all change with the deployment
domain. The engineering goal is not global manifold mastery; it is
identifying the recurring transition types available in the target patch
and provisioning interventions against the head of that local catalogue.

\textbf{Domain transfer is mild but non-trivial.} The mode-rate constant
\(\sigma\) shifts between domains (math \(\approx 1.2\)--\(1.6\); code
\(\approx 0.87\)--\(1.30\); general \(\approx 1.85\)), and
higher-\(\sigma\) domains carry larger catalogues at any given corpus
size --- so they need proportionally more interventions for matched
coverage. The dependence is linear in \(\sigma\) at fixed coverage: a
domain with twice the mode-discovery rate needs roughly twice the
catalogue size. Across the domains measured here, \(\sigma\) does not
vary by orders of magnitude.

\textbf{Per-hard-token vs.~sequence-level targets matter for SLA
design.} The result in Section 3.4 separates two cases that are often
conflated. Per-hard-token residual error rate --- the right metric for
systems that surface or correct errors continuously --- has a polylog
intervention budget. Sequence-level failure probability --- the right
metric for systems where a single error invalidates the output --- is
strictly stricter and requires nearly full-catalogue coverage as \(k\)
grows. Production deployments should design their reliability SLA to
match the actual cost structure: continuous-correction systems can lean
on the polylog result, while one-shot systems should plan for tighter
coverage.

\textbf{Library design as engineering, not asymptotic theory.} A
practical implication of the cluster-selectivity finding (Section 4.2)
is that an intervention library can be assembled module-by-module, in
any order, as long as the modules target distinct failure classes. The
Le (2026) caveat applies only to interventions sharing a prompt channel;
layer-separated interventions (constrained decoding + retrieval +
process supervision + tool call) compose approximately additively in our
data, though the additivity is empirical rather than derived. Library
construction is therefore highly parallelisable: different teams can
ship independent interventions and combine them with minimal integration
risk, modulo the prompt-channel caveat.

\textbf{Capability libraries are coarser than error-class libraries.}
Five named error classes in the math domain (arithmetic, units,
counting, list manipulation, date arithmetic) are all addressed by a
single Python interpreter; three code-error classes (SyntaxError,
NameError, most TypeError) collapse under one execution-feedback loop;
format violations and the structural component of ``missing required
element'' --- together more than 20\% of ErrorAtlas --- collapse under
one constrained-decoder deployment. The practical reliability library is
therefore smaller than the count of named error clusters. A team
budgeting for ``\(\approx 50\) interventions per domain'' can expect to
need fewer than 50 capability-axis interventions while still covering
\(\approx 50\) named clusters. The six axes of Section 4.2.0 are likely
closer to the right unit of engineering accounting than the twelve
clusters of Section 4.2.

\begin{center}\rule{0.5\linewidth}{0.5pt}\end{center}

\subsection{6. Discussion}\label{discussion}

\subsubsection{6.1 By-construction elimination as a boundary
case}\label{by-construction-elimination-as-a-boundary-case}

Seven of the 28 capability-elimination citations in Section 4.2.0 ---
DeepSeek-Prover-V2's Lean proof checker, OpenAI Structured Outputs,
ToolDec, DINGO, XGrammar, XGrammar-2, and LlmFix's static syntax check
--- achieve residual error rate equal to zero \emph{by mathematical
construction}, not statistically. Constrained decoders set
\(P(\text{invalid token}) = 0\) at every generation step, so no invalid
token can be emitted. Proof kernels reject any output that fails
type-checking. Schema validators reject any output missing a required
field. (Acurai's RAG result is excluded from Pattern A here, despite its
100\% RAGTruth number; it is a benchmark-conditional empirical
observation, not a by-construction property of citation-grounded
retrieval in general --- see Section 4.2.0 Pattern B. LlmFix's
\texttt{NameError} repair is also excluded, since its residual is 2.3\%
rather than zero --- Pattern C.)

In this regime, the polylog bound of Proposition 2 ---
\(m \geq \lceil |C|^{1 - \varepsilon / e_{\text{hard}}} \rceil\) --- is
loose. The cluster covered by a by-construction intervention contributes
zero to the residual error, not a small positive quantity. The formal
statement is that for any cluster \(c \in \{1, \dots, m\}\) covered by a
Pattern-A intervention,
\(\Pr(\text{failure} \mid \text{cluster } c) = 0\) identically,
regardless of how many hard-token decisions the model makes inside that
cluster. The Proposition 2 budget formula treats coverage as a
\([0, 1]\) fractional quantity; Pattern A pushes a subset of clusters to
coverage exactly 1.

This is a strengthening result, not a counter-example. The framework's
qualitative claim --- that capability provisioning can erase an error
class --- admits a sharp boundary case where the implication is
mathematical certainty rather than statistical reduction. Crucially,
Pattern A applies precisely to the \emph{structural / verifiable}
classes (format, syntax, schema, invalid proofs), which are also the
head of the heavy-tailed cluster frequency distribution. Capability
provisioning therefore exerts its strongest effect exactly where the
failure-mode catalogue is densest, which compounds the polylog
conclusion of Section 3.4.

\subsubsection{6.2 Re-audit of counter-evidence: relocation, not
dissolution}\label{re-audit-of-counter-evidence-relocation-not-dissolution}

Four prominent papers are routinely cited as evidence that LLM
reliability decays steeply with length, in apparent tension with our
framework: Dziri et al.~(2023), Kuratov et al.~(2024), METR (2025), and
Wan et al.~(2025). To these we add NoCha (Karpinska et al., 2024) for
completeness. A careful re-reading shows that \emph{every one} of these
papers decays over a variable distinct from raw token length \(n\).

Identifying the decay axis is not the same as dissolving the practical
concern. In real tasks where compositional graph size, fact count, or
evidence scope grow with problem length, the framework predicts steep
failure curves; the contribution is to identify which interventions help
(capability provisioning along the actual decay axis) and which do not
(longer context windows, more compute on the \(n\)-axis). The papers
below are not refuted by the framework. Each correctly identifies a
regime in which the framework expects difficulty --- sometimes severe
--- and our analysis below acknowledges those regimes explicitly rather
than relabelling them away.

\textbf{Dziri et al.~(2023), Faith and Fate.} GPT-4 multi-digit
multiplication accuracy drops from 59\% (3-digit) to 4\% (4-digit) in
zero-shot. The authors theorise that ``the probability of incorrect
predictions converges exponentially to \(\approx 1\) for abstract
compositional tasks.'' The decay variable, however, is
\emph{compositional graph size \(N\)}, not raw token length \(n\). The
3×3 multiplication graph has on the order of \(d^2\) partial products
plus carries --- call the resulting node count \(N_3\); the 4×4 graph
similarly has node count \(N_4\), with \(N_4 > N_3\) (Dziri et
al.~report the graph-size grouping but exact node counts depend on the
carry convention, so we read off the order-of-magnitude rather than
commit to specific integers). Multi-digit multiplication is engineered
so that \(k_{\text{hard}} \approx N\) --- every node in the computation
graph is a hard decision. For natural-language tasks, where the harvest
evidence in Section 4.1 gives \(k_{\text{hard}} \ll n\), the Dziri
regime is the \emph{boundary case} in which our framework reduces to
their result. We agree with their empirical decay; we disagree with the
interpretation that it extends to natural-language \(n\).

A complementary observation: a clean \((1 - \varepsilon)^N\) exponential
cannot simultaneously reproduce 59\% at the 3-digit graph size and 4\%
at the 4-digit graph size for any single per-node \(\varepsilon\) ---
the observed drop is locally steeper than per-node-iid exponential. This
is consistent with a finite catalogue of failure modes that gets
exhausted as \(N\) grows, not a smooth uniform per-node failure rate.
\textbf{Where the framework concedes}: in adversarial compositional
tasks engineered so that \(k_{\text{hard}} \approx N\), the framework
reduces to Dziri's regime and does not relieve it; the relocation is
informational, not magical.

\textbf{Kuratov et al.~(2024), BABILong.} The abstract reports
``performance declines sharply with increased reasoning complexity'' and
that models ``effectively utilise only 10--20\% of the context.'' Read
in detail, BABILong varies \emph{two} axes: context length \(n\) (from
0K to 10M tokens) and number of supporting facts \(k\) (from QA1 = 1
fact to QA3 = 3 facts, etc.). The sharp decline is in \(k\), not in
\(n\). For QA1 (single-fact), most models ``perform well up to 4,000
tokens'' --- a plateau, not exponential decay. The famous
RAG-flat-across-length result (60\% on single-fact QA \emph{independent
of context length}) is the cleanest possible demonstration that when the
relevant evidence is in window, length does not matter. Recurrent Memory
Transformers maintaining performance to 50M tokens further confirms that
effective \(k\) is determined by architecture, not by raw \(n\). The
framework's concession is narrow but real: for multi-hop tasks whose
required fact count grows with task complexity, BABILong's sharp
\(k\)-axis decay is exactly what we predict happens, not a
counter-example to it.

\textbf{METR (2025).} Per-model success fits a logistic curve in
\(\log(\text{human-task-duration})\):
\(S(t) = \sigma(\beta \cdot (\log t - \log h))\). Logistic-in-log-length
is \emph{mathematically sublinear} in length itself. The 80\%-horizon
being 4--6× shorter than the 50\%-horizon is a steep within-model cliff
but is by construction compatible with the polylog result of Section 3.4
--- a cliff at a specific capacity threshold is a manifold-transition
signature, exactly the structure of Arbuzov et al.~(2025, §3.2). The
famous exponential --- capability doubling every seven months --- is an
\emph{inter-model} claim about how the horizon \(h\) moves across
generations, orthogonal to within-model decay shape. When a model
improves, its \(h\) shifts right; the curve does not steepen. An honest
qualifier: the within-model logistic-in-log cliff is steep on a
practical scale. Calling it ``sublinear in length'' is technically
correct but engineering-useful only for systems that can be designed to
operate well below the 50\% horizon.

\textbf{Wan et al.~(2025), Fano-style upper bound.} This paper theorises
super-linear information-demand growth and identifies an ``accuracy
cliff'' at capacity overflow: ``when the task's information demand
surpasses the model's output capacity, performance does not degrade
gracefully but instead collapses sharply.'' The behavioural prediction
is a \emph{threshold}, not smooth \((1 - \varepsilon)^n\). A cliff at a
specific capacity threshold is exactly the manifold-transition behaviour
our Section 3 framework predicts at the boundary between covered and
uncovered modes; Wan et al.'s mechanism (capacity overflow) is one
specific cause of the boundary, complementary to ours. \textbf{Where the
framework concedes}: Wan documents a mechanism by which \(|C|\)
effectively exceeds the model's coverage in a single forward pass;
Pattern A interventions (constrained decoding, formal verification) are
the kind of response the framework expects but does not automatically
supply.

\textbf{Karpinska et al.~(2024), NoCha.} GPT-4o achieves 55.8\% pair
accuracy across 1,001 minimally-different true/false claim pairs about
long fictional books (mean length 127K tokens). The paper reports
performance by \emph{evidence scope}, not by context length: 59.8\% on
sentence-level retrieval, 47.6\% on passage-level, 41.6\% on global
reasoning. No per-context-length curve is reported. The decay axis is
the number of evidence pieces that must be integrated --- i.e., a
\(k\)-axis quantity, not an \(n\)-axis one. NoCha is therefore not
counter-evidence to a sublinear-in-\(n\) claim. Evidence-scope decay is
exactly the \(k\)-axis observation the framework predicts cannot be
addressed by scaling raw context length --- it requires retrieval or
process supervision along the actual decay axis. The relocation is
informational, not magical.

\textbf{Unifying observation.} Each of the steep-decay counter-papers we
re-audit decays over a variable other than \(n\): compositional graph
size, fact count, log-time horizon, capacity threshold, or evidence
scope. The apparent rapid decay is, in each case, in a quantity that our
framework already concentrates the action in (\(k_{\text{hard}}\) and
\(|C|\)), not in raw sequence length. This is a relocation, not a
dissolution. These papers are not counter-evidence to the
sublinear-in-\(n\) thesis, but they identify regimes where the framework
predicts steep \(k_{\text{hard}}\)-axis decay in practice. Identifying
the correct intervention axis is the contribution of this re-audit;
promising that those regimes are easy is not. Where \(k_{\text{hard}}\)
grows with task length --- adversarial compositional structure,
multi-hop fact chains, long horizons that force more decisions ---
reliability remains hard; the framework's value is in directing
intervention toward capability provisioning along the actual decay axis
rather than toward context-window or compute-budget expansion that does
not help.

\subsubsection{6.3 Limitations, residual budget, and
falsifiability}\label{limitations-residual-budget-and-falsifiability}

\textbf{Mathematical limitations.} Three are worth flagging. First,
Postulate 1 is empirical, not derived. While it is consistent with
ErrorAtlas, HumanEval, and MWPES at corpus scale, a domain with
genuinely Heaps-power-law mode discovery (canonical \(b \approx 0.5\))
would invalidate the doubly-logarithmic special case of Proposition 2,
although the qualitative polylog result of Appendix A.5 would still
hold. Second, the coverage form in Section 3.3 is an empirical best-fit
choice among several plausible candidates; readers should not interpret
the close numerical match against two anchor points as theoretical
confirmation. Third, the cluster-orthogonality assumption that
underwrites Proposition 2's composition step is only approximate; the Le
(2026) caveat on prompt-channel interference is real and would tighten
the bound in shared-channel settings.

\textbf{Domain narrowness.} Postulate 1's empirical anchor rests on
three taxonomies (ErrorAtlas, HumanEval categorisation, MWPES-300K) all
published in 2025--2026 and covering general, code, and math domains.
The framework is untested on agentic workflows, long-running scientific
reasoning, multi-turn tool use over millions of tokens, and production
codebases. Those are precisely the regimes where \(|C|\) may grow faster
than logarithmically --- and where the practical reliability question
matters most. The polylog conclusion of Section 3.4 should be read as a
domain-conditional claim until subsample-discovery measurements
(Appendix A.5 lists the falsifiability test) are available in those
regimes.

\textbf{Single-point calibration of \(\sigma\).} The estimate
\(\sigma \approx 1.85\) is a calibration against the ErrorAtlas category
count at one corpus scale, not a fit to a discovery curve. A
subsample-vs-distinct-modes plot has not been published for ErrorAtlas
(or for any other LLM error taxonomy at the time of writing); Appendix
A.5 lists this as the explicit empirical test that would either tighten
the postulate or move the analysis to the Heaps variant.

\textbf{Identifiability of \(\beta\).} \(\beta\) is a latent
stratification parameter; we report no empirical range because
failure-mode atlases measure
\(\Pr(\text{category} \mid \text{error occurred})\), which does not
identify \(\Pr(\text{hard} \mid \text{key-token})\). Future work
measuring hard-token opportunities directly --- e.g., by ablation of
model confidence at predicted decision boundaries, or by comparing
confidence distributions over sampled key-token positions --- would
close this gap.

\textbf{Taxonomy granularity and the L2--L3 gap.} The framework operates
on L2 (taxonomy categories) as a proxy for L3 (latent modes). \(|C|\)
depends on the taxonomer's resolution choices: splitting ``reasoning
error'' into 200 subtypes inflates \(|C|\) without changing the
underlying mode catalogue. We have no independent measurement of how
faithful current taxonomies are as L3 proxies, and we treat this as a
residual epistemic risk rather than a solved problem.

\textbf{Goalpost-relocation caveat.} By concentrating the practical
action in \(k_{\text{hard}}\) and \(|C|\) rather than raw \(n\), the
framework \emph{relocates} the difficulty of long-context reliability
rather than resolving it. Domains where \(k_{\text{hard}}\) grows with
task length --- adversarial compositional structure, multi-hop chains,
long agent horizons --- remain hard. The framework's claim is to
identify the axis of intervention (capability provisioning along the
actual decay variable), not to dissolve the underlying engineering
problem. Practical reliability work in those regimes still requires the
hard interventions; the framework's value is in saying which
interventions are on-axis and which are off.

\textbf{Patch-shift between deployments.} When a model moves from one
application domain to another --- say, from general-purpose chat to
clinical decision support --- the four domain-indexed quantities
\(A_D\), \(\sigma_D\), \(\beta_D\), and \(|C_D|\) all change. A library
calibrated against one patch will under-cover the next patch's mode
catalogue. The framework predicts this transfer failure structurally
(heavy-tailed and bounded per patch) but says nothing about the
\emph{content} of \(C_{D'}\) in advance. The operationally useful
response is to re-measure inside the new patch and re-fit the head of
the local distribution. The cleanest follow-up evidence would be
per-domain mode-discovery measurement --- subsamples of size \(t\)
versus distinct discovered modes, the falsifiability test of Appendix
A.5 --- but no such curve has been published at the time of writing. The
framework's prediction is that such curves should rise log- or
Heaps-style inside the patch and asymptote to a domain-specific
\(|C_D|\).

\textbf{The irreducible-semantic residual.} Of the 17 ErrorAtlas
categories, 13 are addressed under Patterns A, B, or C by one of the six
capability axes of Section 4.2.0; four are not. The mapping below makes
this audit-able through the paper's own 12-cluster scheme of Section 4.2
(with Category I split into its structural and semantic sub-clusters per
§4.2's reframing, yielding 13 distinct addressable sub-classes):

{\def\LTcaptype{none} % do not increment counter
\begin{longtable}[]{@{}
  >{\raggedright\arraybackslash}p{(\linewidth - 6\tabcolsep) * \real{0.2500}}
  >{\raggedright\arraybackslash}p{(\linewidth - 6\tabcolsep) * \real{0.2500}}
  >{\raggedright\arraybackslash}p{(\linewidth - 6\tabcolsep) * \real{0.2500}}
  >{\raggedright\arraybackslash}p{(\linewidth - 6\tabcolsep) * \real{0.2500}}@{}}
\toprule\noalign{}
\begin{minipage}[b]{\linewidth}\raggedright
Cluster (this paper)
\end{minipage} & \begin{minipage}[b]{\linewidth}\raggedright
Covering axis
\end{minipage} & \begin{minipage}[b]{\linewidth}\raggedright
Pattern
\end{minipage} & \begin{minipage}[b]{\linewidth}\raggedright
Notes
\end{minipage} \\
\midrule\noalign{}
\endhead
\bottomrule\noalign{}
\endlastfoot
A. Arithmetic / computation & Arithmetic & B & PAL, PoT, ToRA \\
B. Unit conversion & Arithmetic & B & folded into A --- Python handles
units \\
C. Counting / enumeration & Arithmetic & (gap → C) & smallest residual
gap \\
D. Format / schema violation & Format/Structure & A & constrained
decoding \\
E. Code logic / type / null & Code Exec & A/B & REPL feedback;
SyntaxError is Pattern A \\
F. Multi-hop decomposition drift & Knowledge/RAG & C & entity-grounded
rewriter \\
G. Reasoning step errors & Verification & B & process reward \\
H. Specification misinterpretation & Verification (weak) & C &
clarification loops; weakest fit \\
I.a Structural missing element & Format/Structure & A & subsumed by D's
constrained decoders \\
I.b Semantic missing element & Verification + clarification & C &
subsumed by G/H \\
J. Hallucination / factual & Knowledge/RAG & B/C & RAG, citation
grounding \\
K. Inappropriate refusal & Verification & C & preference optimisation \\
L. Tool / API usage & Format/Structure + Verification & B & structured
uncertainty \\
\end{longtable}
}

That accounts for 13 sub-classes mapped to one of the six axes. The
remaining four ErrorAtlas categories --- \emph{(1)} a residual of
inappropriate refusal beyond what preference optimisation reaches,
\emph{(2)} specification misinterpretation not closed by clarification
loops, \emph{(3)} reasoning bottlenecks at the level of problem
decomposition, and \emph{(4)} a small ``user wanted something
different'' semantic remainder --- admit no clean
capability-provisioning fix in the current literature. These are
\emph{irreducible-semantic residuals}: classes where the failure is in
choosing what to do, not in executing it. Proposition 2's prediction of
\(O(\log)\) rather than \(O(0)\) residual reliability reflects exactly
this irreducible core. The framework does not claim 100\% coverage of
all failures; it claims polylog-bounded \emph{capability-eliminable}
failures, with the named semantic residual as the floor.

\textbf{A falsifiability test the framework passes.} DebugBench (Tian et
al., 2024) explicitly catalogues which error classes execution-feedback
can and cannot repair. It works well for syntax and reference errors ---
quoted: \emph{``SyntaxError can be fixed: 99.74\% on HumanEval, 100.00\%
on MBPP.''} It does not work for logic errors --- quoted: \emph{``the
feedback information is unhelpful for logic errors\ldots{} may even
cause disruptions.''} This is the framework's predicted failure mode
(execution feedback is a Pattern-A/B capability for structural-class
errors and a Pattern-C-at-best for semantic-class errors) and a
falsifiability test it passes. If capability provisioning bled into the
semantic residual --- if execution feedback claimed to fix logic errors
--- the framework's selectivity claim would be wrong. The empirical
result is that capability provisioning correctly identifies both what it
can erase and what it cannot.

\textbf{Net assessment.} The framework correctly characterises both the
success regime and the failure regime; the residuals are honest
limitations, not silent failures. A future extension addressing the
irreducible-semantic residual --- through richer process supervision,
ensemble-based verification, or higher-order capability layers --- would
tighten the bound but is not required for the polylog conclusion to hold
over the capability-eliminable fraction.

\begin{center}\rule{0.5\linewidth}{0.5pt}\end{center}

\subsection{7. Conclusion}\label{conclusion}

LLM reliability is often framed as an asymptotic scaling problem: as
outputs get longer, do errors compound to inevitable failure? In our
previous work (Arbuzov et al., 2025) we argued that this framing rests
on a false uniformity assumption --- errors concentrate at
\(\sim 5{-}10\%\) of tokens, not all of them. This paper takes the next
step: within that sparse set, errors are not only concentrated but also
\emph{repetitive}. They cluster into a finite catalogue of recurring
failure modes whose size, under the empirical postulate of Section 3.2,
grows logarithmically --- and under the more conservative Heaps
alternative, as a small power --- of the number of observed failures.

The conditional consequence (Proposition 2) is that a sufficient
per-hard-decision intervention budget scales polylogarithmically in
sequence length within a fixed domain patch under the log
active-mode-exposure assumption, and becomes a domain-constant once the
patch ceiling \(|C_D|\) is reached. The optimistic special case
(Postulate 1 tight, \(k \sim \log n\)) gives a doubly-logarithmic rate;
the conservative Heaps variant of Appendix A.5 gives \((\log n)^b\) with
small \(b\); the symbolic-rates form of Appendix A.5 shows the polylog
conclusion is stable across this family. Available evidence is
consistent with libraries on the order of tens of interventions covering
the head of the per-hard-decision failure distribution in many fixed
domains. The exact number is patch-indexed and should be estimated from
local discovery and rank-coverage curves, not assumed from cross-domain
priors. Sequence-level reliability targets are strictly tighter and
approach full-catalogue coverage as \(k\) grows; practitioners should
pick the regime that matches their cost structure rather than reading
the per-decision result as the production SLA.

The same logarithmic discovery postulate carries an inverse
interpretation (Corollary 1, Appendix A.6): the logarithmic envelope
cannot deliver linearly more distinct tail modes without exponentially
more observed hard-failure events. Combined with head-heavy failure
mass, this explains why generic frontier post-training can face
diminishing reliability returns in open-ended deployment settings ---
the high-mass head is discovered early, while the tail becomes
increasingly expensive to find and contributes less marginal
residual-error reduction. The framework therefore does not say frontier
scaling is useless. It says frontier-only reliability is economically
misaligned with fixed deployment patches: once \(D\) is known, local
adaptation, tools, validators, retrieval, constrained decoding, and
process supervision can target recurring failure mass directly, where
generic post-training would have to rediscover the same local repairs
indirectly across an open-ended task universe.

This reframes the long-context reliability question from ``can we bound
the growth of \(n\)-token error?'' to ``have we catalogued enough
failure modes?'' The latter question is finite and addressable. Whether
it generalises to agentic, scientific, and long-horizon regimes where
the published taxonomies do not yet reach is the open empirical question
this paper invites. In the domains where taxonomies have been measured
--- general, code, math --- the architecture of LLM errors is compact
and the engineering problem is small. Whether the same holds elsewhere
is a falsifiable empirical test, and we have named it.

Two empirical extensions stand out. Direct measurement of the mode-rate
constant \(\sigma\) on new domains --- agentic, scientific,
code-in-production --- would test the cross-domain generality of
Postulate 1 and surface outlier domains where the catalogue may grow
faster than logarithmically. A fuller sequence-level validation ---
measuring \(S_{\text{base}}\), \(\beta\), and active-mode exposure
directly in deployed systems --- would turn the sequence-level analogue
of §3.4 from a formal bound into an operational SLA estimator.

The conceptual follow-on goes in a different direction. This paper names
the engineering object: a patch-local failure catalogue and the
intervention budget that covers its head. It does not say how production
systems actually accumulate and govern that library over time. That
question --- how the localhost scaffold (instructions, tools, retrieval,
memory, orchestration, governance) learns to cover the local failure
topology --- is the subject of Part 3 of this series, \emph{Frontier and
Localhost}. Reliability engineering, in that frame, is governance of the
scaffold rather than scaling of the weights.

\begin{center}\rule{0.5\linewidth}{0.5pt}\end{center}

\subsection{Appendix A. Formal Proofs, Derivations, and Sensitivity
Analysis}\label{appendix-a.-formal-proofs-derivations-and-sensitivity-analysis}

\subsubsection{A.1 Proof of Proposition 1: No Universal Finite
Intervention
Dictionary}\label{a.1-proof-of-proposition-1-no-universal-finite-intervention-dictionary}

Fix a residual-error tolerance \(\varepsilon\). Say that an intervention
\emph{covers} a failure mode if it reduces the residual error of that
mode below \(\varepsilon\). Coverage is mode-level: an intervention that
covers a mode covers every failure event in that mode. The two
conditions ``\(D\) is intervention-unbounded'' and
``\(|C_D^{\varepsilon}| = \infty\)'' are equivalent under mode-level
coverage; an unbounded witnessing sequence is constructed by taking one
representative per mode, and an infinite catalogue forces the existence
of such a sequence.

Let \(D\) be a domain. Call \(D\) \emph{intervention-unbounded} if it
contains an infinite sequence of failures \(f_1, f_2, f_3, \ldots\) such
that each new \(f_j\) is not covered by any finite intervention
dictionary that covers all earlier failures
\(\{f_1, \ldots, f_{j-1}\}\).

\textbf{Step 1: assume the opposite.} Suppose, for contradiction, that
\(C_D^{\varepsilon}\) is finite. Then there are only finitely many
intervention-distinguishable failure modes in \(D\). Write them as
\(C_D^{\varepsilon} = \{c_1, \ldots, c_M\}\) for some finite \(M\).

\textbf{Step 2: what finiteness means.} If there are only \(M\)
intervention-distinguishable modes, then after all \(M\) modes have
appeared in the sequence, every later failure must belong to one of the
already-seen modes.

\textbf{Step 3: same mode means same intervention class.} Modes are
defined at intervention resolution \(\varepsilon\). If a later failure
belongs to the same mode as an earlier failure, then the intervention
dictionary that covers the earlier representative of that mode also
covers the later failure below residual tolerance \(\varepsilon\).

\textbf{Step 4: contradiction.} Intervention-unboundedness says exactly
the opposite: each new \(f_j\) is not covered by any finite dictionary
that covers \(\{f_1, \ldots, f_{j-1}\}\). Therefore \(f_j\) cannot
belong to any earlier intervention mode, so each \(f_j\) introduces a
new intervention-distinguishable mode. The sequence
\(f_1, f_2, f_3, \ldots\) induces infinitely many such modes,
contradicting the assumption that \(C_D^{\varepsilon}\) is finite. Hence
\(|C_D^{\varepsilon}| = \infty\).

\textbf{Step 5: no finite dictionary covers every mode of the domain.}
Suppose, again for contradiction, that some finite intervention
dictionary \(\mathcal{I}\) covers every mode of \(D\). Then
\(\mathcal{I}\) covers every finite prefix \(\{f_1, \ldots, f_{j-1}\}\)
for every \(j\). By intervention-unboundedness, any dictionary covering
that prefix fails to cover \(f_j\). Therefore \(\mathcal{I}\) does not
cover \(f_j\), contradicting the assumption that \(\mathcal{I}\) covers
every mode of \(D\). \(\square\)

\subsubsection{A.2 Derivation of Proposition 2: Patch-Local Sufficient
Intervention
Budget}\label{a.2-derivation-of-proposition-2-patch-local-sufficient-intervention-budget}

\textbf{Conditions used.} Proposition 2 gives a sufficient,
model-implied budget under the following assumptions:

\begin{enumerate}
\def\labelenumi{\arabic{enumi}.}
\tightlist
\item
  \textbf{Coverage model.} The cumulative hard-error mass covered by the
  top \(m\) modes is approximated by
  \(F_{\log}(m; |C|) = \min(1,\, \ln m / \ln |C|)\) on the declared
  domain \(m \geq 2\), \(|C| \geq 2\). Throughout this appendix, \(F\)
  refers to \(F_{\log}\) unless otherwise stated.
\item
  \textbf{Non-trivial, attainable target.} The residual target satisfies
  \(\varepsilon \in (0, e_{\text{hard}})\). If
  \(\varepsilon \geq e_{\text{hard}}\) no intervention is needed; if
  \(\varepsilon \leq 0\) the target is unattainable unless all residual
  hard-token error is eliminated.
\item
  \textbf{Patch-local catalogue.} The result applies only after a
  deployment patch \(D\) has been fixed and its reachable catalogue is
  modelled as finite or effectively capped.
\item
  \textbf{Sequence-length claim.} The doubly-logarithmic rate further
  requires Assumption 2 and \(k(n) = \Theta(\log n)\).
\end{enumerate}

\textbf{Step 1: define the target.} Let \(e_{\text{hard}}\) be the
baseline hard-token error rate and
\(\varepsilon \in (0, e_{\text{hard}})\) the target after intervention.

\textbf{Step 2: define the effective catalogue.} For per-sequence
scaling, set \(C_{\text{eff}} = C_{\text{active},D}(n)\). For full
deployment-library budgeting, set \(C_{\text{eff}} = C_D\).

\textbf{Step 3: cumulative coverage.} Let the ranked local failure modes
have hard-error masses
\(p_1 \geq p_2 \geq \cdots \geq p_{|C_{\text{eff}}|}\) with
\(\sum_i p_i = 1\). A library covering the top \(m\) modes removes
cumulative hard-error mass
\(F(m; |C_{\text{eff}}|) = \sum_{i=1}^{m} p_i\), approximated by the
log-coverage form.

\textbf{Step 4: residual after intervention.} The uncovered fraction is
\(1 - F(m; |C_{\text{eff}}|)\), so the residual per-hard-token error
rate is
\(e_{\text{res}}(m) = (1 - F(m; |C_{\text{eff}}|))\, e_{\text{hard}}\).

\textbf{Step 5: impose the target.} Requiring
\(e_{\text{res}}(m) \leq \varepsilon\) and dividing by
\(e_{\text{hard}} > 0\) gives
\(F(m; |C_{\text{eff}}|) \geq 1 - \varepsilon / e_{\text{hard}}\).

\textbf{Step 6: substitute the log-coverage form.} Since
\(\varepsilon < e_{\text{hard}}\), the required coverage
\(1 - \varepsilon/e_{\text{hard}} \in (0, 1)\), so the cap in \(F\) is
inactive before saturation. Using \(F = \ln m / \ln |C_{\text{eff}}|\),

\[
\frac{\ln m}{\ln |C_{\text{eff}}|} \;\geq\; 1 - \frac{\varepsilon}{e_{\text{hard}}}.
\]

\textbf{Step 7: solve for \(m\).} Multiplying by
\(\ln |C_{\text{eff}}| > 0\) and exponentiating gives
\(m \geq |C_{\text{eff}}|^{1 - \varepsilon / e_{\text{hard}}}\). Taking
the ceiling (since \(m\) is integer-valued):

\[
m \;\geq\; \left\lceil |C_{\text{eff}}|^{1 - \varepsilon / e_{\text{hard}}} \right\rceil.
\]

\textbf{Step 8: sequence-length rate.} For per-sequence scaling,
\(C_{\text{eff}} = C_{\text{active},D}(n)\). Assumption 2 bounds
\(|C_{\text{active},D}(n)| \leq \min(A'_D + \sigma'_D \ln h(n),\, |C_D|)\).
In the pre-cap regime, while \(A'_D + \sigma'_D \ln h(n) < |C_D|\), the
ceiling has not been reached and
\(|C_{\text{active},D}(n)| = O(\ln h(n))\). With \(h(n) = \beta k(n)\)
and \(k(n) = \Theta(\log n)\), \(\ln h(n) = \Theta(\log\log n)\), hence
\(|C_{\text{active},D}(n)| = O(\log\log n)\). Substituting into the
budget gives

\[
m \;=\; O\!\bigl((\log\log n)^{1 - \varepsilon / e_{\text{hard}}}\bigr).
\]

\textbf{Step 9: cap regime.} Once active-mode discovery saturates the
patch catalogue, \(|C_{\text{active},D}(n)| = |C_D|\), so
\(m \geq \lceil |C_D|^{1 - \varepsilon / e_{\text{hard}}} \rceil\), no
longer depending on \(n\). \(\square\)

\subsubsection{A.3 Sequence-Level Reliability
Derivation}\label{a.3-sequence-level-reliability-derivation}

Proposition 2 gives a per-hard-token residual target. Production systems
often care about a stricter target: the probability that the entire
sequence is correct.

Starting from the composed reliability of §3.1, and writing the
post-intervention hard-token rate as \(e_{\text{res}}\):

\[
P(\text{correct}) \;=\; (1 - e_{\text{res}})^{\beta k}\,(1 - e_{\text{easy}})^{(1 - \beta) k}\,(1 - e_{\text{non}})^{n - k}.
\]

Define the non-hard-token survival factor

\[
S_{\text{base}} \;=\; (1 - e_{\text{easy}})^{(1 - \beta) k}\,(1 - e_{\text{non}})^{n - k},
\]

so that
\(P(\text{correct}) = (1 - e_{\text{res}})^{\beta k}\, S_{\text{base}}\).
Requiring \(P(\text{correct}) \geq 1 - \varepsilon_{\text{seq}}\)
becomes

\[
(1 - e_{\text{res}})^{\beta k} \;\geq\; \frac{1 - \varepsilon_{\text{seq}}}{S_{\text{base}}}.
\]

Three regimes follow.

\textbf{Regime (i): non-hard-token errors already violate the target.}
If \(S_{\text{base}} < 1 - \varepsilon_{\text{seq}}\), even setting
\(e_{\text{res}} = 0\) cannot meet the target, because the maximum
possible survival after eliminating all hard-token failures is only
\(S_{\text{base}}\). Hard-token interventions alone cannot meet the
sequence-level SLA. This is the back-door route by which the original
exponential-in-\(n\) concern can re-enter, and it should be diagnosed
before any catalogue-budgeting exercise.

\textbf{Regime (ii): hard-token residual error determines feasibility.}
If \(S_{\text{base}} \geq 1 - \varepsilon_{\text{seq}}\), the target may
be feasible. Taking the \((1/(\beta k))\)-th power and isolating
\(e_{\text{res}}\),

\[
e_{\text{res}} \;\leq\; 1 - \left(\frac{1 - \varepsilon_{\text{seq}}}{S_{\text{base}}}\right)^{1 / (\beta k)}
\;=:\; \tau_{\text{seq}}.
\]

Applying Proposition 2 with \(\varepsilon\) replaced by
\(\tau_{\text{seq}}\) gives the sufficient library size
\(m \geq \lceil |C_{\text{eff}}|^{1 - \tau_{\text{seq}} / e_{\text{hard}}} \rceil\).
As the sequence target becomes stricter (\(\varepsilon_{\text{seq}}\)
shrinks), \(\tau_{\text{seq}} \to 0\) and the exponent
\(1 - \tau_{\text{seq}}/e_{\text{hard}} \to 1\), so
\(m \to |C_{\text{eff}}|\). Strict one-shot sequence-level reliability
pushes the system toward full-catalogue coverage.

\textbf{Regime (iii): baseline hard-token error is already acceptable.}
If \(\tau_{\text{seq}} \geq e_{\text{hard}}\), the baseline hard-token
rate already satisfies the sequence target (since
\(e_{\text{res}}(0) = e_{\text{hard}}\)).

\textbf{The additive shortcut as a heuristic.} Taking logs and using
\(\log(1-x) \approx -x\) for small \(x\) gives the familiar

\[
-\log P(\text{correct}) \;\approx\; \beta k\, e_{\text{res}}(m) + (1 - \beta) k\, e_{\text{easy}} + (n - k)\, e_{\text{non}},
\]

and the hard-token term dominates \emph{only} when
\((n-k)\, e_{\text{non}} = o(\beta k\, e_{\text{res}}(m))\) ---
equivalently \(e_{\text{non}} = o(\beta \alpha\, e_{\text{res}}(m))\)
with \(\alpha = k/n\). Under that shortcut,
\(e_{\text{res}}(m) \lesssim \varepsilon_{\text{seq}}/(\beta k)\) and
the budget reduces to
\(m \geq \lceil |C|^{1 - \varepsilon_{\text{seq}}/(\beta k\, e_{\text{hard}})} \rceil\).
We retain this as a production heuristic; the multiplicative
\(S_{\text{base}}/\tau_{\text{seq}}\) treatment above is the literal
statement, and the shortcut is valid only inside Regime (ii) and only
when the non-key aggregate is genuinely small relative to
\(\beta k\, e_{\text{res}}(m)\).

\textbf{Conclusion.} Per-hard-token reliability is easier than
sequence-level reliability. A library that gives a large reduction in
residual hard-token error may still be insufficient for one-shot
sequence-level guarantees when many hard decisions occur in a single
output. This is why the main paper treats the ``tens of interventions''
rule of thumb as a per-hard-decision planning prior, not a one-shot
sequence-level SLA.

\subsubsection{A.4 Why these are propositions rather than unconditional
theorems}\label{a.4-why-these-are-propositions-rather-than-unconditional-theorems}

Proposition 1 is a definitional impossibility result: once
intervention-unboundedness is assumed, an infinite
intervention-resolution catalogue follows. Its role is not to prove that
every unbounded domain necessarily has infinite failure modes, but to
show that open-ended domains cannot be assumed to admit finite
dictionaries.

Proposition 2 is conditional engineering math. It does not prove that
LLM failures universally obey logarithmic mode discovery. It proves that
\emph{if} a bounded patch has a finite or effectively capped reachable
catalogue, \emph{and if} cumulative intervention coverage follows the
stated head-heavy form, \emph{then} a sufficient per-hard-decision
intervention budget grows slowly and becomes constant after catalogue
saturation. It does not prove that the true minimal intervention library
has the same scaling.

The empirical burden therefore lies not in the algebra but in measuring,
for each deployment patch, the local discovery curve
\(C_{\text{seen},D}(T)\), the per-sequence activation
\(C_{\text{active},D}(n)\), the cumulative coverage \(F(m; |C_D|)\), the
hard-token fraction \(\beta_D\), and the baseline hard-token rate
\(e_{\text{hard}}\). This is why the paper frames the results as a
reliability-engineering scaffold rather than as universal theorems about
LLM behaviour.

\subsubsection{A.5 Heaps power-law variant and cluster-count
sensitivity}\label{a.5-heaps-power-law-variant-and-cluster-count-sensitivity}

A reader who prefers to derive cluster-count growth from standard Heaps'
law rather than from Postulate 1 obtains a qualitatively similar result.
Let \(|C|(k_{\text{hard}}) \approx K \cdot k_{\text{hard}}^{b}\) with
\(b \in (0, 1)\). Canonical fits give \(b \approx 0.5\) for
natural-language vocabularies; failure-mode taxonomies plateau much more
sharply (ErrorAtlas stabilises at \(|C| = 17\) across \(10^4{+}\)
failures), implying a smaller effective \(b \approx 0.2\) in our
setting. Composing with \(k = \Theta(\log n)\), we get
\(|C| = O((\log n)^{b})\), and Proposition 2 becomes

\[
m = O\!\left((\log n)^{b \cdot (1 - \varepsilon / e_{\text{hard}})}\right).
\]

This is still polylogarithmic in \(n\) for any \(b \in (0, 1)\) and any
\(\varepsilon < e_{\text{hard}}\). The paper's qualitative claim ---
that the intervention budget grows polylogarithmically in sequence
length --- survives either choice of cluster-count law.

\textbf{Symbolic-form sensitivity across candidate laws.} The polylog
conclusion of Proposition 2 depends on which cluster-count law one
accepts. The available evidence (§3.2) is consistent with multiple
candidates: no subsample-discovery curve has been published for any LLM
failure-mode taxonomy at the time of writing. We therefore report
symbolic rates rather than fitted constants. Assume
\(h(n) = \beta k(n)\). The candidate active-catalogue laws then imply:

\begin{itemize}
\tightlist
\item
  \textbf{Logarithmic:} \(|C_{\text{active},D}(n)| = O(\log h(n))\). If
  \(k(n) = \Theta(\log n)\), then
  \(m = O\!\bigl((\log\log n)^{1 - \varepsilon / e_{\text{hard}}}\bigr)\).
\item
  \textbf{Heaps:} \(|C_{\text{active},D}(n)| = O\!\bigl(h(n)^b\bigr)\)
  with \(b \in (0, 1)\). If \(k(n) = \Theta(\log n)\), then
  \(m = O\!\bigl((\log n)^{b\,(1 - \varepsilon / e_{\text{hard}})}\bigr)\).
\item
  \textbf{Saturating:} \(|C_{\text{active},D}(n)| \leq |C_D|\). Then
  \(m = O\!\bigl(|C_D|^{1 - \varepsilon / e_{\text{hard}}}\bigr)\),
  constant in \(n\) once the patch ceiling is reached.
\end{itemize}

The qualitative conclusion is robust \textbf{under the
\(k(n) = \Theta(\log n)\) regime adopted throughout}: \(m\) then grows
more slowly than any positive power of \(n\) under every candidate law
above, and the directional claim (``a small library covers the head of
the failure distribution in the per-hard-token regime'') survives. Only
the exponent shifts: doubly-logarithmic under logarithmic discovery,
\((\log n)^b\) with small \(b\) under Heaps, constant in the cap regime.
The doubly-logarithmic rate is the optimistic special case. If \(k(n)\)
grows as a positive power of \(n\), the Heaps variant inherits that
power and the polylog-in-\(n\) language fails along that axis --- the
framework's intervention prescription still applies, but its
asymptotic-rate framing does not.

Numerical constants require a measured discovery curve
\(C_{\text{seen},D}(T)\) or \(C_{\text{active},D}(n)\). Existing
taxonomies provide endpoint category counts at a single corpus scale
(e.g., ErrorAtlas at \(|C| = 17\) for \(\approx 10^4\) failures), not
discovery curves. We therefore report rates rather than fitted
constants, and treat the subsample-discovery curve as the explicit
empirical test that would either tighten the postulate or fall back to
the Heaps variant. Until that measurement exists, the framework's
headline rate should be read as \emph{polylogarithmic in the pre-cap
regime, domain-constant in the cap regime} --- with the specific
exponent flagged as a falsifiability test rather than a fitted
prediction.

\subsubsection{A.6 Inverse Discovery
Cost}\label{a.6-inverse-discovery-cost}

The body Corollary 1 (§3.2) inverts the upper-bound discovery postulate
into a sample-budget lower bound on novel-mode discovery. This appendix
gives the algebra, the numerical anchors, the tightness assumption that
converts the lower bound into an approximate inverse cost, sensitivity
to the Heaps cluster-count alternative of §A.5, the saturation regime,
and a separate mode-mediated gain sub-corollary that connects discovery
cost to broad capability proxies.

\textbf{Setup.} Let \(q(T) = |C_{\text{seen},D}(T)|\) be the number of
distinct failure modes discovered in patch \(D\) after \(T\) observed
hard-failure events. Postulate 1 (§3.2) is an upper bound, \[
q(T) \leq A_D + \sigma_D \ln T, \qquad \sigma_D > 0,
\] on the discovered catalogue under the logarithmic-envelope
assumption.

\textbf{Inversion (lower bound on \(T\)).} The upper-bound postulate is
monotone in \(T\). Taking the inverse direction yields a \emph{lower
bound} on \(T\) required for the envelope to permit \(q\) discovered
modes: if \(q > A_D\) distinct modes have been discovered after \(T\)
events, then \[
T \;\geq\; \exp\!\left(\frac{q - A_D}{\sigma_D}\right).
\] This is the rigorous direction of the corollary. The upper envelope
cannot reach \(q\) until the sample budget has grown exponentially in
\(q\).

\textbf{Tightness assumption.} The lower bound above is unconditional
under Postulate 1. The stronger reading --- that observed hard failures
at scale \(T(q) \approx \exp((q - A_D)/\sigma_D)\) actually deliver
\(q\) discovered modes, and that each additional \(\Delta q\) modes
raises the sample budget by approximately \(\exp(\Delta q/\sigma_D)\)
--- requires the additional assumption that the empirical discovery
curve is approximately \emph{tight} against the envelope at the relevant
corpus scales. Without that assumption, the appendix gives a lower bound
on \(T\) only. With it, \[
\frac{T(q + \Delta q)}{T(q)} \;\approx\; \exp\!\left(\frac{\Delta q}{\sigma_D}\right).
\]

\textbf{Numerical anchors under tightness.} At the conservative
calibration \(\sigma_D \approx 1.85\) (§3.2):

\begin{itemize}
\tightlist
\item
  \(\exp(5/1.85) \approx 14.9\) --- five extra modes need roughly
  \(15\times\) more observed hard failures.
\item
  \(\exp(10/1.85) \approx 222\) --- ten extra modes need roughly
  \(220\times\) more.
\end{itemize}

These are tightness-conditional anchors, not unconditional consequences
of Postulate 1. The subsample-discovery measurement of §3.2 is precisely
the test of whether tightness holds.

\textbf{What this is and is not.} The corollary describes \emph{new
distinct-mode discovery}. Ordinary failures inside already-discovered
modes may remain common and cheap to observe; the exponential cost lives
only on the \emph{category-novelty} axis. Conflating ordinary failure
rate with novel-mode arrival rate would over-claim the result.

\textbf{Heaps alternative.} Under the Heaps cluster-count law of A.5,
\(q(T) = K T^b\) with \(b \in (0,1)\), the corresponding inverse cost is
polynomial rather than exponential: \[
T(q) = (q / K)^{1/b}.
\] The exponential inverse-cost reading is specific to the logarithmic
upper envelope; the broader qualitative claim --- that tail discovery
has diminishing returns --- survives under any concave discovery curve.
Sensitivity is therefore: exponential under logarithmic-and-tight,
polynomial under Heaps, undefined past the patch ceiling.

\textbf{Saturation regime.} Once \(q(T) \leq |C_D|\) has been saturated,
discovery stops; the inversion above applies only in the pre-cap regime.
Inside saturated patches the corollary's multiplicative-cost reading is
vacuous because no novel modes remain to discover, which is itself a
property of the patch, not a failure of the corollary.

\textbf{Mode-mediated capability gain (a second sub-corollary).} Suppose
broad capability or reliability gain \(G\) inside the patch is
approximately linear in the number of useful discovered modes,
\(G(q) = G_0 + \gamma q\) for some \(\gamma > 0\). Composing with the
logarithmic envelope at tightness, \[
G(T) \;\approx\; G_0 + \gamma A_D + \gamma \sigma_D \ln T,
\] so \(G\) grows logarithmically in observed hard-failure exposure
under tightness. Inverting, \[
T(G) \;\approx\; \exp\!\left(\frac{G - G_0 - \gamma A_D}{\gamma \sigma_D}\right).
\] A linear gain in mode-mediated broad reliability therefore
corresponds to exponential growth in observed hard-failure exposure
under the postulate at tightness. We phrase \(G\) as a
\emph{mode-mediated capability/reliability} proxy, not as
``intelligence'': the corollary does not say frontier scaling is useless
or that intelligence requires exponential data in any general sense. The
narrower claim is that for fixed deployment reliability where
improvement is mediated by discovering new useful modes, generic
open-domain training pays a heavy data tax relative to direct
patch-local measurement and intervention.

\textbf{Engineering reading.} The corollary explains, without invoking
new mechanisms, two empirical signals: (a) why generic post-training
shows diminishing reliability returns once a domain's head modes are
covered, and (b) why patch-local measurement combined with targeted
tools, retrieval, validators, constrained decoding, and process
supervision often outperforms more frontier-scale data on the deployment
SLA. Frontier scaling and patch-local engineering solve different
problems: scaling improves the substrate; patch-local engineering
removes recurring deployment failure mass.

\begin{center}\rule{0.5\linewidth}{0.5pt}\end{center}

\subsection{Appendix B. Patch Evidence
Detail}\label{appendix-b.-patch-evidence-detail}

The body's claim that domain patches cap the engineering problem rests
on a triangulation of indirect evidence. None of the items below
\emph{measures} the patch ceiling \(|C_D|\) directly; they support the
weaker claim that model behaviour is strongly domain-dependent and that
operational neighbourhoods occupy bounded regions of the model's
behavioural space.

{\def\LTcaptype{none} % do not increment counter
\begin{longtable}[]{@{}
  >{\raggedright\arraybackslash}p{(\linewidth - 4\tabcolsep) * \real{0.3333}}
  >{\raggedright\arraybackslash}p{(\linewidth - 4\tabcolsep) * \real{0.3333}}
  >{\raggedright\arraybackslash}p{(\linewidth - 4\tabcolsep) * \real{0.3333}}@{}}
\toprule\noalign{}
\begin{minipage}[b]{\linewidth}\raggedright
Evidence type
\end{minipage} & \begin{minipage}[b]{\linewidth}\raggedright
What it supports
\end{minipage} & \begin{minipage}[b]{\linewidth}\raggedright
What it does \emph{not} prove
\end{minipage} \\
\midrule\noalign{}
\endhead
\bottomrule\noalign{}
\endlastfoot
Low intrinsic-dimensional structure in representations (Park et al.,
2024; Li \& Sarwate, 2025) & Domains occupy localised structure in model
representation space & Does not measure failure-catalogue size
\(\lvert C_D \rvert\) \\
Cross-domain accuracy spreads on a single base model (Wang Y. et al.,
2024) & Same model behaves differently across domains; performance is
patch-dependent & Does not imply finite failure modes \\
Long-tail / popularity thresholds in knowledge tasks (Mallen et al.,
2023; Kandpal et al., 2023) & Coverage depends on domain frequency and
training exposure & Does not prove logarithmic mode discovery \\
Patch-specific intervention performance (across §4.2) & Local tools,
schemas, and capability libraries change residual error structure & Does
not imply universal cross-patch transfer of the same intervention
library \\
\end{longtable}
}

The purpose of this evidence is motivational. It supports patch-indexing
of \(\sigma_D\), \(A_D\), \(\beta_D\), \(|C_D|\), but the finite (or
effectively capped) reachable catalogue remains an empirical modelling
assumption to be measured per deployment, not derived from any of the
rows above.

\begin{center}\rule{0.5\linewidth}{0.5pt}\end{center}

\subsection{Appendix C. Empirical Calibration
Detail}\label{appendix-c.-empirical-calibration-detail}

Three published LLM error taxonomies anchor the mode-rate parameter
\(\sigma\) in the body's \(\sigma \in [0.87, 1.85]\) range. The simple
calibration uses \(|C| \approx A + \sigma \ln T\) with \(A = 0\); this
is a deliberately conservative readout that ignores any positive
intercept and uses endpoint counts rather than discovery curves.

{\def\LTcaptype{none} % do not increment counter
\begin{longtable}[]{@{}
  >{\raggedright\arraybackslash}p{(\linewidth - 8\tabcolsep) * \real{0.1667}}
  >{\raggedleft\arraybackslash}p{(\linewidth - 8\tabcolsep) * \real{0.2222}}
  >{\raggedleft\arraybackslash}p{(\linewidth - 8\tabcolsep) * \real{0.2222}}
  >{\raggedleft\arraybackslash}p{(\linewidth - 8\tabcolsep) * \real{0.2222}}
  >{\raggedright\arraybackslash}p{(\linewidth - 8\tabcolsep) * \real{0.1667}}@{}}
\toprule\noalign{}
\begin{minipage}[b]{\linewidth}\raggedright
Source
\end{minipage} & \begin{minipage}[b]{\linewidth}\raggedleft
Approx. observed failures \(T\)
\end{minipage} & \begin{minipage}[b]{\linewidth}\raggedleft
Named categories \(\lvert C \rvert\)
\end{minipage} & \begin{minipage}[b]{\linewidth}\raggedleft
Implied \(\sigma\) at \(A=0\)
\end{minipage} & \begin{minipage}[b]{\linewidth}\raggedright
Role
\end{minipage} \\
\midrule\noalign{}
\endhead
\bottomrule\noalign{}
\endlastfoot
ErrorAtlas (Ashury-Tahan et al., 2026) & \(\gtrsim 10^4\) & 17 &
\(\approx 1.85\) & Conservative cross-domain anchor (used as planning
value) \\
HumanEval categorisation (Wen et al., 2024) & comparable scale &
\(8\)--\(12\) & \(\approx 0.87\)--\(1.30\) & Code-domain anchor \\
MWPES-300K (Sun et al., 2025) & \(\approx 3 \times 10^5\) &
\(15\)--\(20\) & \(\approx 1.2\)--\(1.6\) & Math-domain anchor (largest
corpus) \\
\end{longtable}
}

These are endpoint counts, not discovery curves: they tell us \emph{how
many} named categories a taxonomer assigned at a single corpus scale,
not how the count grew with \(T\). They do not prove logarithmic mode
discovery. Their role in this paper is twofold: they motivate the
empirical postulate of §3.2, and they identify the explicit empirical
test that would either tighten the postulate or move the analysis to the
Heaps variant of Appendix A.5: repeatedly subsample failures from a
fixed deployment patch \(D\) and plot discovered modes
\(|C_{\text{seen},D}(T)|\) against \(T\).

\begin{center}\rule{0.5\linewidth}{0.5pt}\end{center}

\subsection{References}\label{references}

Arbuzov, M. L., Bei, S., Dong, Z., Kalaev, D., \& Shvets, A. A. (2025).
Beyond exponential decay: Rethinking error accumulation in large
language models. \emph{arXiv preprint arXiv:2505.24187v2 {[}cs.CL{]}},
06 May 2026. CC BY 4.0.

Asai, A., Wu, Z., Wang, Y., Sil, A., \& Hajishirzi, H. (2024). Self-RAG:
Learning to retrieve, generate, and critique through self-reflection.
\emph{International Conference on Learning Representations (ICLR) 2024}.
arXiv:2310.11511.

Ashury-Tahan, S., Mai, Y., Bandel, E., Shmueli-Scheuer, M., \& Choshen,
L. (2026). ErrorMap and ErrorAtlas: Charting the failure landscape of
large language models. \emph{arXiv preprint arXiv:2601.15812}.

Cemri, M., Pan, M. Z., Yang, S., Agrawal, L. A., Chopra, B., Tiwari, R.,
Keutzer, K., Parameswaran, A., Klein, D., Ramchandran, K., Zaharia, M.,
Gonzalez, J. E., \& Stoica, I. (2025). Why do multi-agent LLM systems
fail? \emph{arXiv preprint arXiv:2503.13657}. (Introduces the MAST
taxonomy.)

Chen, W., Ma, X., Wang, X., \& Cohen, W. W. (2023). Program of thoughts
prompting: Disentangling computation from reasoning for numerical
reasoning tasks. \emph{Transactions on Machine Learning Research (TMLR)
2023}. arXiv:2211.12588.

Cheng, K., Sun, Q., Chu, Y., Xu, F., Li, Y., Zhang, J., \& Wu, Z.
(2024). SeeClick: Harnessing GUI grounding for advanced visual GUI
agents. \emph{Proceedings of ACL 2024}. arXiv:2401.10935.

Costello, C., Guo, S., Goldie, A., \& Mirhoseini, A. (2025). Think,
prune, train, improve: Scaling reasoning without scaling models.
\emph{arXiv preprint arXiv:2504.18116}.

Dai, C., et al.~(2025). Capture the key in reasoning to enhance CoT
distillation generalization. \emph{Proceedings of ACL 2025}.
arXiv:2405.19737. (Earlier arXiv title: ``Beyond Imitation: Learning Key
Reasoning Steps from Dual Chain-of-Thoughts in Reasoning Distillation''.
Referred to in body as EDIT.)

Dantart, A. (2026). Reliability by design: Quantifying and eliminating
fabrication risk in LLMs. From generative to consultative AI: A
comparative analysis in the legal domain and lessons for high-stakes
knowledge bases. \emph{arXiv preprint arXiv:2601.15476}.

Dong, Y., Ruan, C. F., Cai, Y., Lai, R., Xu, Z., Zhao, Y., \& Chen, T.
(2025). XGrammar: Flexible and efficient structured generation engine
for large language models. \emph{Proceedings of MLSys 2025}.
arXiv:2411.15100. (Cited as 2024 in body.)

Li, L., Dong, Y., Wang, G., Xu, Z., Jiang, A., \& Chen, T. (2026).
XGrammar-2: Efficient dynamic structured generation engine for agentic
LLMs. \emph{arXiv preprint arXiv:2601.04426}.

Dziri, N., Lu, X., Sclar, M., Li, X. L., Jiang, L., Lin, B. Y., West,
P., Bhagavatula, C., Le Bras, R., Hwang, J. D., et al.~(2023). Faith and
fate: Limits of transformers on compositionality. \emph{Advances in
Neural Information Processing Systems 36}. arXiv:2305.18654.

Fang, L., Wang, Y., Liu, Z., Zhang, C., Jegelka, S., Gao, J., Ding, B.,
\& Wang, Y. (2025). What is wrong with perplexity for long-context
language modeling? \emph{International Conference on Learning
Representations (ICLR) 2025}. arXiv:2410.23771. (Cited as 2024 in some
earlier drafts; 2025 is the ICLR publication year.)

Gao, L., Madaan, A., Zhou, S., Alon, U., Liu, P., Yang, Y., Callan, J.,
\& Neubig, G. (2023). PAL: Program-aided language models.
\emph{Proceedings of ICML 2023}. arXiv:2211.10435. (Cited as 2022 in
body.)

Gao, T., Yen, H., Yu, J., \& Chen, D. (2023). Enabling large language
models to generate text with citations. \emph{Proceedings of EMNLP
2023}. arXiv:2305.14627. (Introduces the ALCE benchmark.)

Geng, S., et al.~(2025). JSONSchemaBench: A rigorous benchmark of
structured outputs for language models. \emph{arXiv preprint
arXiv:2501.10868}.

Goodell, A. J., Chu, S. N., Rouholiman, D., \& Chu, L. F. (2025). Large
language model agents can use tools to perform clinical calculations.
\emph{npj Digital Medicine}, 8(1), 163. DOI: 10.1038/s41746-025-01475-8.

Gou, B., Wang, R., Zheng, B., Xie, Y., Chang, C., Shu, Y., Sun, H., \&
Su, Y. (2025). Navigating the digital world as humans do: Universal
visual grounding for GUI agents. \emph{International Conference on
Learning Representations (ICLR) 2025}. arXiv:2410.05243. (Introduces the
UGround model; ICLR 2025 Oral.)

Gou, Z., Shao, Z., Gong, Y., Shen, Y., Yang, Y., Huang, M., Duan, N., \&
Chen, W. (2024). ToRA: A tool-integrated reasoning agent for
mathematical problem solving. \emph{International Conference on Learning
Representations (ICLR) 2024}. arXiv:2309.17452. (Cited as 2023 in body.)

Gou, Z., Shao, Z., Gong, Y., Shen, Y., Yang, Y., Duan, N., \& Chen, W.
(2024). CRITIC: Large language models can self-correct with
tool-interactive critiquing. \emph{International Conference on Learning
Representations (ICLR) 2024}. arXiv:2305.11738.

Guo, D., Liu, J., Fan, Z., He, Z., Li, H., Li, Y., Wang, Y., \& Fung, Y.
R. (2025). Mathematical proof as a litmus test: Revealing failure modes
of advanced large reasoning models. \emph{arXiv preprint
arXiv:2506.17114}. (Introduces the RFMDataset.)

Hsieh, C.-P., Sun, S., Kriman, S., Acharya, S., Rekesh, D., Jia, F.,
Zhang, Y., \& Ginsburg, B. (2024). RULER: What's the real context size
of your long-context language models? \emph{Conference on Language
Modeling (COLM) 2024}. arXiv:2404.06654.

Huang, C., Zhu, G., Wang, X., Luo, Y., Ge, G., Chen, H., Yi, D., \&
Wang, J. (2024). Recurrent context compression: Efficiently expanding
the context window of LLM. \emph{arXiv preprint arXiv:2406.06110}.

Huang, D., Bu, Q., Zhang, J. M., Luck, M., \& Cui, H. (2023).
AgentCoder: Multi-agent-based code generation with iterative testing and
optimisation. \emph{arXiv preprint arXiv:2312.13010}.

Jiang, H., Wu, Q., Luo, X., Li, D., Lin, C.-Y., Yang, Y., \& Qiu, L.
(2024). LongLLMLingua: Accelerating and enhancing LLMs in long context
scenarios via prompt compression. \emph{Proceedings of ACL 2024}.
arXiv:2310.06839.

Kandpal, N., Deng, H., Roberts, A., Wallace, E., \& Raffel, C. (2023).
Large language models struggle to learn long-tail knowledge.
\emph{Proceedings of ICML 2023}. arXiv:2211.08411.

Karaman, B. K., Zabir, I., Benhaim, A., Chaudhary, V., Sabuncu, M. R.,
\& Song, X. (2024). POROver: Improving safety and reducing overrefusal
in large language models with overgeneration and preference
optimization. \emph{arXiv preprint arXiv:2410.12999}.

Karpinska, M., Thai, K., Lo, K., Goyal, T., \& Iyyer, M. (2024). One
thousand and one pairs: A ``novel'' challenge for long-context language
models. \emph{Proceedings of EMNLP 2024}. arXiv:2406.16264.

Kuratov, Y., Bulatov, A., Anokhin, P., Rodkin, I., Sorokin, D., Sorokin,
A., \& Burtsev, M. (2024). BABILong: Testing the limits of LLMs with
long context reasoning-in-a-haystack. \emph{NeurIPS 2024 Datasets and
Benchmarks Track}. arXiv:2406.10149.

Kwa, T., et al.~(2025). Measuring AI ability to complete long software
tasks. \emph{Advances in Neural Information Processing Systems (NeurIPS)
2025}. arXiv:2503.14499. (METR technical report.)

Le, Y. (2026). Schema key wording as an instruction channel in
structured generation under constrained decoding. \emph{arXiv preprint
arXiv:2604.14862}.

LeCun, Y. (2023). Auto-regressive LLMs are doomed. \emph{Tweet, March
26, 2023.} \url{https://x.com/ylecun/status/1640122342570336267}.

Li, X., \& Sarwate, A. D. (2025). Unraveling the localized latents:
Learning stratified manifold structures in LLM embedding space with
sparse mixture-of-experts. \emph{arXiv preprint arXiv:2502.13577}.

Li, Y., et al.~(2022). Competition-level code generation with AlphaCode.
\emph{Science}, 378(6624), 1092--1097. DOI: 10.1126/science.abq1158.

Lightman, H., Kosaraju, V., Burda, Y., Edwards, H., Baker, B., Lee, T.,
Leike, J., Schulman, J., Sutskever, I., \& Cobbe, K. (2024). Let's
verify step by step. \emph{International Conference on Learning
Representations (ICLR) 2024}. arXiv:2305.20050. (Cited as 2023 in body.)

Liu, F., Eisenschlos, J. M., Piccinno, F., Krichene, S., Pang, C., Lee,
K., Joshi, M., Chen, W., Collier, N., \& Altun, Y. (2023). DePlot:
One-shot visual language reasoning by plot-to-table translation.
\emph{Findings of ACL 2023}. arXiv:2212.10505.

Lu, Y., Yang, J., Shen, Y., \& Awadallah, A. (2024). OmniParser for pure
vision based GUI agent. \emph{arXiv preprint arXiv:2408.00203}.

Mallen, A., Asai, A., Zhong, V., Das, R., Khashabi, D., \& Hajishirzi,
H. (2023). When not to trust language models: Investigating
effectiveness of parametric and non-parametric memories.
\emph{Proceedings of ACL 2023}. arXiv:2212.10511. (Introduces the PopQA
benchmark.)

Manning, C. D., Raghavan, P., \& Schütze, H. (2008). \emph{Introduction
to information retrieval}. Cambridge University Press. (Section on
Heaps' law:
\url{https://nlp.stanford.edu/IR-book/html/htmledition/heaps-law-estimating-the-number-of-terms-1.html}.)

Niwa, A., \& Iso, H. (2024). AmbigNLG: Addressing task ambiguity in
instruction for NLG. \emph{Proceedings of EMNLP 2024}. arXiv:2402.17717.

OpenAI. (2024). Introducing structured outputs in the API. \emph{OpenAI
Engineering Blog, August 6, 2024.}
\url{https://openai.com/index/introducing-structured-outputs-in-the-api/}.

OpenAI. (2025). GPT-5 system card. \emph{Published August 7, 2025.}
\url{https://cdn.openai.com/gpt-5-system-card.pdf}.

Pang, J., Ye, F., Wong, D. F., He, X., Chen, W., \& Wang, L. (2024).
Anchor-based large language models. \emph{Findings of ACL 2024}.
arXiv:2402.07616.

Park, K., Choe, Y. J., \& Veitch, V. (2024). The linear representation
hypothesis and the geometry of large language models. \emph{Proceedings
of ICML 2024}. arXiv:2311.03658.

Patel, K., Surendira, S., George, J., \& Kapale, S. (2026). The Six
Sigma Agent: Achieving enterprise-grade reliability in LLM systems
through consensus-driven decomposed execution. \emph{arXiv preprint
arXiv:2601.22290}.

Press, O., Zhang, M., Min, S., Schmidt, L., Smith, N. A., \& Lewis, M.
(2023). Measuring and narrowing the compositionality gap in language
models. \emph{Findings of EMNLP 2023}. arXiv:2210.03350.

Ren, Z. Z., et al.~(2025). DeepSeek-Prover-V2: Advancing formal
mathematical reasoning via reinforcement learning for subgoal
decomposition. \emph{arXiv preprint arXiv:2504.21801}.

Schaeffer, R., et al.~(2025). How do large language monkeys get their
power (laws)? \emph{International Conference on Machine Learning (ICML)
2025}. arXiv:2502.17578.

Schick, T., Dwivedi-Yu, J., Dessì, R., Raileanu, R., Lomeli, M., Hambro,
E., Zettlemoyer, L., Cancedda, N., \& Scialom, T. (2023). Toolformer:
Language models can teach themselves to use tools. \emph{Advances in
Neural Information Processing Systems 36}. arXiv:2302.04761.

Shang, Y., Li, Y., Zhao, K., Ma, L., Liu, J., Xu, F., \& Li, Y. (2024).
AgentSquare: Automatic LLM agent search in modular design space.
\emph{arXiv preprint arXiv:2410.06153}.

Shi, Y., Wang, S., Wan, C., Wang, M., \& Gu, X. (2024). From code to
correctness: Closing the last mile of code generation with hierarchical
debugging. \emph{arXiv preprint arXiv:2410.01215}. (Introduces the
MGDebugger method.)

Shinn, N., Cassano, F., Berman, E., Gopinath, A., Narasimhan, K., \&
Yao, S. (2023). Reflexion: Language agents with verbal reinforcement
learning. \emph{Advances in Neural Information Processing Systems 36}.
arXiv:2303.11366.

Sun, Y., Yin, Z., Huang, X., Qiu, X., \& Zhao, H. (2025). Error
classification of large language models on math word problems: A
dynamically adaptive framework. \emph{arXiv preprint arXiv:2501.15581}.
(Introduces the MWPES-300K dataset.)

Suresh, T., et al.~(2025). DINGO: Constrained inference for diffusion
LLMs. \emph{arXiv preprint arXiv:2505.23061}.

Suri, M., Mathur, P., Lipka, N., Dernoncourt, F., Rossi, R. A., \&
Manocha, D. (2025). Structured uncertainty guided clarification for LLM
agents. \emph{arXiv preprint arXiv:2511.08798}. (Introduces the
SAGE-Agent method.)

Tian, R., Ye, Y., Qin, Y., Cong, X., Lin, Y., Liu, Z., \& Sun, M.
(2024). DebugBench: Evaluating debugging capability of large language
models. \emph{Findings of ACL 2024}. arXiv:2401.04621.

Wada, A., et al.~(2025). Retrieval-augmented generation elevates local
LLM quality in radiology contrast media consultation. \emph{npj Digital
Medicine}, 8(1), 395. DOI: 10.1038/s41746-025-01802-z.

Wan, K., et al.~(2026). A Fano-style accuracy upper bound for LLM
single-pass reasoning in multi-hop QA. \emph{International Conference on
Learning Representations (ICLR) 2026}. arXiv:2509.21199. (Cited as 2025
in body.)

Wang, B., et al.~(2025). From scores to steps: Diagnosing and improving
LLM performance in evidence-based medical calculations.
\emph{Proceedings of EMNLP 2025}. arXiv:2509.16584. (Introduces the
MedRaC benchmark.)

Wang, D. Y.-B., Shen, Z., Mishra, S. S., Xu, Z., Teng, Y., \& Ding, H.
(2025). SLOT: Structuring the output of large language models.
\emph{arXiv preprint arXiv:2505.04016}. (SLOT = Structured LLM Output
Transformer.)

Wang, M., et al.~(2024). Leave no document behind: Benchmarking
long-context LLMs with extended multi-doc QA. \emph{Proceedings of EMNLP
2024}. arXiv:2406.17419.

Wang, P., Li, L., Shao, Z., Xu, R. X., Dai, D., Li, Y., Chen, D., Wu,
Y., \& Sui, Z. (2024). Math-Shepherd: Verify and reinforce LLMs
step-by-step without human annotations. \emph{Proceedings of ACL 2024}.
arXiv:2312.08935. (Cited as 2023 in body.)

Wang, X., Wei, J., Schuurmans, D., Le, Q., Chi, E., Narang, S.,
Chowdhery, A., \& Zhou, D. (2023). Self-consistency improves chain of
thought reasoning in language models. \emph{International Conference on
Learning Representations (ICLR) 2023}. arXiv:2203.11171.

Wang, Y., et al.~(2024). MMLU-Pro: A more robust and challenging
multi-task language understanding benchmark. \emph{NeurIPS 2024 Datasets
and Benchmarks Track (Spotlight)}. arXiv:2406.01574.

Wen, H., et al.~(2024). Fixing function-level code generation errors for
foundation large language models. \emph{arXiv preprint
arXiv:2409.00676}. (Introduces the LlmFix method. Body cites as Wen et
al.; an earlier draft mis-attributed to ``Zhang et al.'')

Wood, M. C., \& Forbes, A. A. (2024). 100\% elimination of
hallucinations on RAGTruth for GPT-4 and GPT-3.5 Turbo. \emph{arXiv
preprint arXiv:2412.05223}. (Introduces the Acurai method.)

Xie, T., et al.~(2025). Scaling computer-use grounding via user
interface decomposition and synthesis. \emph{Advances in Neural
Information Processing Systems (NeurIPS) 2025}. arXiv:2505.13227.
(NeurIPS 2025 Spotlight. Introduces the Jedi grounding dataset and
OSWorld-G benchmark.)

Yao, S., Shinn, N., Razavi, P., \& Narasimhan, K. (2024).
\(\tau\)-bench: A benchmark for tool-agent-user interaction in
real-world domains. \emph{arXiv preprint arXiv:2406.12045}.

Zakka, C., Shad, R., Chaurasia, A., Dalal, A. R., Kim, J. L., Moor, M.,
Fong, R., Phillips, C., Alexander, K., Ashley, E., Boyd, J., Boyd, K.,
Hirsch, K., Langlotz, C., Lee, R., Melia, J., Nelson, J., Sallam, K.,
Tullis, S., Vogelsong, M. A., Cunningham, J. P., \& Hiesinger, W.
(2024). Almanac --- Retrieval-augmented language models for clinical
medicine. \emph{NEJM AI}, 1(2), AIoa2300068. DOI: 10.1056/AIoa2300068.

Zhang, K., et al.~(2023). Don't fine-tune, decode: Syntax error-free
tool use via constrained decoding. \emph{arXiv preprint
arXiv:2310.07075}. (Introduces the ToolDec method.)

Zhang, M., Meng, Z., \& Collier, N. (2026). Failure modes in multi-hop
QA: The weakest link effect and the recognition bottleneck. \emph{arXiv
preprint arXiv:2601.12499}.

Zhou, Y., Liu, H., Chen, Z., Tian, Y., \& Chen, B. (2025). GSM-Infinite:
How do your LLMs behave over infinitely increasing context length and
reasoning complexity? \emph{arXiv preprint arXiv:2502.05252}.

Zhu, R., Liu, X., Sun, Z., Wang, Y., \& Hu, W. (2025). Mitigating
lost-in-retrieval problems in retrieval-augmented multi-hop question
answering. \emph{Proceedings of ACL 2025}. arXiv:2502.14245.

\end{document}
